% =============================================================
% RAPPORT DE PROJET DE FIN D'ÉTUDES — AEGISWEB
% Plateforme de Gouvernance et de Confiance pour Agents IA Web
%
% Compilateur : pdfLaTeX
% Encodage    : UTF-8
% Cible       : Overleaf (pdfLaTeX) ou compilation locale
% =============================================================

\documentclass[12pt,a4paper,oneside]{report}

% ── Encodage & Langue ─────────────────────────────────────────
\usepackage[utf8]{inputenc}
\usepackage[T1]{fontenc}
\usepackage[french]{babel}

% ── Polices ───────────────────────────────────────────────────
\usepackage{mathptmx}
\usepackage[scaled=0.90]{helvet}
\usepackage{courier}

% ── Mise en page ──────────────────────────────────────────────
\usepackage[
  left        = 3cm,
  right       = 2cm,
  top         = 2.5cm,
  bottom      = 2.5cm,
  headheight  = 22pt,
  headsep     = 10pt,
  footskip    = 20pt
]{geometry}

\usepackage{setspace}
\onehalfspacing

\setlength{\parskip}{3pt plus 1pt minus 1pt}
\setlength{\parindent}{0pt}

% ── Couleurs (palette AegisWeb) ───────────────────────────────
\usepackage{xcolor}
\definecolor{aegisblue}{RGB}{10, 54, 120}
\definecolor{aegisaccent}{RGB}{0, 112, 200}
\definecolor{aegislightblue}{RGB}{232, 242, 255}
\definecolor{aegisgray}{RGB}{100, 100, 100}
\definecolor{codebg}{RGB}{250, 250, 252}
\definecolor{codeframe}{RGB}{200, 210, 220}
\definecolor{successgreen}{RGB}{39, 174, 96}
\definecolor{dangerred}{RGB}{192, 57, 43}
\definecolor{warningorange}{RGB}{211, 84, 0}
\definecolor{rowgray}{RGB}{245, 247, 250}

% ── En-tête & Pied de page ────────────────────────────────────
\usepackage{fancyhdr}
\usepackage{graphicx}

\pagestyle{fancy}
\fancyhf{}
\renewcommand{\headrulewidth}{0.5pt}
\renewcommand{\footrulewidth}{0.5pt}
\renewcommand{\headrule}{%
  \hbox to\headwidth{%
    \color{aegisblue}\leaders\hrule height\headrulewidth\hfill}}
\renewcommand{\footrule}{%
  \hbox to\headwidth{%
    \color{aegisblue}\leaders\hrule height\footrulewidth\hfill}}

\fancyhead[L]{\small
\includegraphics[height=14pt]{assets/aegisweb_standalone_icon.png}}
\fancyhead[C]{\small\textcolor{aegisgray}{\textit{\nouppercase{\leftmark}}}}
\fancyhead[R]{\small\textcolor{aegisblue}{\textsc{AegisWeb}}}
\fancyfoot[L]{\small\textcolor{aegisgray}{Année Universitaire 2025/2026}}
\fancyfoot[C]{\small\textcolor{aegisblue}{\thepage}}
\fancyfoot[R]{\small\textcolor{aegisgray}{[Prénom \textsc{Nom}]}}

\fancypagestyle{plain}{%
  \fancyhf{}
  \renewcommand{\headrulewidth}{0pt}
  \renewcommand{\footrulewidth}{0.5pt}
  \fancyfoot[C]{\small\textcolor{aegisblue}{\thepage}}
}

% ── Style des chapitres ───────────────────────────────────────
\usepackage{titlesec}

\titleformat{\chapter}[display]
  {\normalfont\fontsize{16}{20}\bfseries\color{aegisblue}\centering}
  {\MakeUppercase{\chaptertitlename}~\thechapter}
  {8pt}
  {\fontsize{16}{20}\bfseries\MakeUppercase}
  [\vspace{4pt}{\color{aegisaccent}\titlerule[1.5pt]}]
\titlespacing*{\chapter}{0pt}{10pt}{18pt}

\titleformat{\section}
  {\normalfont\fontsize{14}{18}\bfseries\color{aegisblue}}
  {\thesection}{0.8em}{}
\titlespacing*{\section}{0pt}{12pt}{5pt}

\titleformat{\subsection}
  {\normalfont\fontsize{12}{16}\bfseries\color{aegisblue}}
  {\thesubsection}{0.8em}{}
\titlespacing*{\subsection}{0pt}{9pt}{4pt}

\titleformat{\subsubsection}
  {\normalfont\fontsize{12}{16}\bfseries\itshape\color{aegisblue}}
  {\thesubsubsection}{0.8em}{}
\titlespacing*{\subsubsection}{0pt}{7pt}{3pt}

% ── Table des matières ────────────────────────────────────────
\usepackage{tocloft}
\renewcommand{\cftchapdotsep}{\cftdotsep}
\renewcommand{\cftchapfont}{\bfseries\color{aegisblue}}
\renewcommand{\cftchappagefont}{\bfseries\color{aegisblue}}
\setlength{\cftbeforechapskip}{6pt}
\setlength{\cftbeforesecskip}{2pt}

% ── Tableaux ──────────────────────────────────────────────────
\usepackage{booktabs}
\usepackage{tabularx}
\usepackage{multirow}
\usepackage{array}
\usepackage{longtable}
\usepackage{colortbl}
\usepackage{makecell}
\setlength{\LTcapwidth}{\textwidth}
\renewcommand{\arraystretch}{1.28}
\newcolumntype{C}[1]{>{\centering\arraybackslash}p{#1}}
\newcolumntype{L}[1]{>{\raggedright\arraybackslash}p{#1}}

% ── Figures & Légendes ────────────────────────────────────────
\usepackage{caption}
\usepackage{subcaption}
\usepackage{float}
\captionsetup{
  font          = small,
  labelfont     = {bf,color=aegisblue},
  justification = centering,
  skip          = 5pt
}
\captionsetup[table]{
  position  = above,
  labelfont = {bf,color=aegisblue},
  skip      = 4pt
}

% ── Code source ───────────────────────────────────────────────
\usepackage{listings}

\lstdefinestyle{tsstyle}{
  backgroundcolor   = \color{codebg},
  commentstyle      = \color{green!55!black}\itshape,
  keywordstyle      = \color{aegisblue}\bfseries,
  numberstyle       = \tiny\color{gray},
  stringstyle       = \color{red!60!black},
  basicstyle        = \ttfamily\footnotesize,
  breaklines        = true,
  captionpos        = b,
  numbers           = left,
  numbersep         = 7pt,
  tabsize           = 2,
  frame             = single,
  rulecolor         = \color{codeframe},
  xleftmargin       = 15pt,
  framexleftmargin  = 15pt,
  aboveskip         = 8pt,
  belowskip         = 6pt,
  morekeywords      = {async,await,import,export,const,interface,type,
                       enum,implements,extends,readonly,private,public,
                       protected,static,return,from,true,false,null,
                       undefined,Promise,Record,Partial},
}

% ── Boîtes colorées ───────────────────────────────────────────
\usepackage{tcolorbox}
\tcbuselibrary{skins,breakable}

\tcbset{
  aegiscubox/.style={
    enhanced, breakable,
    boxrule   = 0.8pt,
    colframe  = aegisblue,
    colback   = aegislightblue,
    coltitle  = white,
    fonttitle = \bfseries\small,
    attach boxed title to top left={yshift=-2mm, xshift=4mm},
    boxed title style={
      colback  = aegisblue, colframe = aegisblue,
      boxrule  = 0pt, sharp corners
    },
    top = 7pt, bottom = 7pt,
    before upper = {\setlength{\parskip}{3pt}},
  }
}

% ── Divers ────────────────────────────────────────────────────
\usepackage{enumitem}
\setlist{itemsep=3pt, topsep=4pt, parsep=0pt}
\usepackage{amssymb}
\usepackage{pifont}
\usepackage{mdframed}
\usepackage{lscape}
\usepackage{pgfgantt}
\usepackage{microtype}

\usepackage{chngcntr}
\counterwithin{figure}{chapter}
\counterwithin{table}{chapter}

\usepackage[hidelinks]{hyperref}
\hypersetup{
  colorlinks = true,
  linkcolor  = aegisblue,
  urlcolor   = aegisaccent,
  citecolor  = aegisblue,
  pdftitle   = {Rapport PFE -- AegisWeb -- Gouvernance Agents IA},
  pdfauthor  = {[Prenom NOM]},
  pdfsubject = {Conception et Developpement de AegisWeb}
}

\newmdenv[
  topline         = false,
  bottomline      = false,
  rightline       = false,
  linewidth       = 3pt,
  linecolor       = aegisaccent,
  innerleftmargin = 10pt,
  innerrightmargin= 6pt,
  innertopmargin  = 6pt,
  innerbottommargin = 6pt,
  backgroundcolor = blue!3!white,
  skipabove       = 8pt,
  skipbelow       = 8pt,
]{chapintro}

\newcommand{\diagramespace}[2][5cm]{%
  \begin{center}%
    \setlength{\fboxsep}{10pt}%
    \setlength{\fboxrule}{1pt}%
    \fcolorbox{aegisaccent!60}{gray!4}{%
      \begin{minipage}[c][#1][c]{0.87\textwidth}%
        \centering
        {\color{aegisgray}\large$\square$\enspace}%
        {\color{aegisgray}\textit{\textbf{Espace réservé :} #2}}%
      \end{minipage}%
    }%
  \end{center}%
}

\newcommand{\coche}{\textcolor{successgreen}{\ding{51}}}
\newcommand{\croix}{\textcolor{dangerred}{\ding{55}}}
\newcommand{\attention}{\textcolor{warningorange}{\ding{115}}}

\newcommand{\conclusionchap}[1]{%
  \vspace{10pt}%
  \begin{chapintro}%
    \textbf{Conclusion du chapitre~:} #1%
  \end{chapintro}%
  \vspace{4pt}%
}

% =============================================================
%                      DOCUMENT
% =============================================================
\begin{document}

\pagenumbering{roman}

% =============================================================
% PAGE DE GARDE
% =============================================================
\begin{titlepage}
\thispagestyle{empty}
\centering

\vspace*{0.2cm}
{\large\bfseries RÉPUBLIQUE TUNISIENNE}\\[3pt]
{\small Ministère de l'Enseignement Supérieur et de la Recherche Scientifique}

\vspace{0.4cm}
{\color{aegisgray}\rule{\textwidth}{0.4pt}}
\vspace{0.4cm}

\begin{minipage}[c]{0.38\textwidth}
  \centering
  \fbox{\parbox[c][2.8cm][c]{3.5cm}{%
    \centering\small\textcolor{aegisgray}{Logo\\[4pt]Établissement}}}
\end{minipage}%
\hfill%
\begin{minipage}[c]{0.38\textwidth}
  \centering
  
\includegraphics[height=2.8cm]{assets/aegisweb_primary_logo.png}
\end{minipage}

\vspace{0.5cm}
{\Large\bfseries [NOM DE L'UNIVERSITÉ / ÉCOLE]}\\[4pt]
{\normalsize Département d'Informatique}

\vspace{0.5cm}
{\color{aegisblue}\rule{\textwidth}{2pt}}\\[3pt]
{\Large\bfseries\color{aegisblue} RAPPORT DE PROJET DE FIN D'ÉTUDES}\\[3pt]
{\color{aegisblue}\rule{\textwidth}{2pt}}

\vspace{0.5cm}
{\normalsize\itshape Présenté en vue de l'obtention du}\\[4pt]
{\bfseries [Diplôme National de Licence Appliquée /
Diplôme d'Ingénieur / Mastère Professionnel]}\\[4pt]
{\bfseries Spécialité : Génie Logiciel et Systèmes Distribués}

\vspace{0.7cm}

\begin{mdframed}[
  linecolor       = aegisblue,
  linewidth       = 1.5pt,
  innertopmargin  = 12pt,
  innerbottommargin = 12pt,
  innerleftmargin = 14pt,
  innerrightmargin= 14pt,
  backgroundcolor = aegislightblue
]
\centering
{\Large\bfseries\color{aegisblue}
  Conception et Développement de \textit{AegisWeb}}\\[6pt]
{\normalsize\bfseries\color{aegisblue}
  Une Plateforme de Gouvernance et de Confiance\\
  pour les Agents d'Intelligence Artificielle Agissant sur le Web}
\end{mdframed}

\vspace{0.8cm}

\begin{minipage}[t]{0.45\textwidth}
  \raggedright
  \textbf{Réalisé par :}\\[5pt]
  [Prénom \textsc{Nom} Étudiant]
\end{minipage}%
\hfill%
\begin{minipage}[t]{0.45\textwidth}
  \raggedleft
  \textbf{Encadré par :}\\[5pt]
  [Titre. Prénom \textsc{Nom} Encadreur]\\[8pt]
  \textbf{Encadreur entreprise :}\\[5pt]
  [Titre. Prénom \textsc{Nom}]
\end{minipage}

\vspace{0.8cm}
{\color{aegisgray}\rule{\textwidth}{0.4pt}}\\[5pt]
\textbf{Organisme d'accueil :}~[Nom de l'entreprise / Projet personnel]
\hfill
\textbf{Année Universitaire :}~2025/2026\\[5pt]
{\color{aegisblue}\rule{\textwidth}{1pt}}

\end{titlepage}

% =============================================================
% DÉDICACES
% =============================================================
\clearpage
\chapter*{Dédicaces}
\thispagestyle{plain}
\addcontentsline{toc}{chapter}{Dédicaces}
\markboth{Dédicaces}{Dédicaces}

\vspace*{2.5cm}
\begin{center}
\begin{minipage}{0.78\textwidth}
\centering
\textit{%
À mes chers parents,\\[5pt]
pour leur amour inconditionnel, leur soutien sans failles\\
et les innombrables sacrifices qu'ils ont consentis\\
pour me permettre d'atteindre cet objectif.\\[1cm]
À ma famille et à mes proches,\\[5pt]
dont la présence bienveillante et les encouragements m'ont\\
accompagné tout au long de ce parcours.\\[1cm]
À mes camarades et amis,\\[5pt]
avec qui j'ai partagé les joies, les défis\\
et les moments inoubliables de ces années d'études.\\[1.2cm]
\textbf{Je vous dédie ce modeste travail.}%
}
\end{minipage}
\end{center}

% =============================================================
% REMERCIEMENTS
% =============================================================
\clearpage
\chapter*{Remerciements}
\thispagestyle{plain}
\addcontentsline{toc}{chapter}{Remerciements}
\markboth{Remerciements}{Remerciements}

Nous tenons à exprimer notre profonde gratitude à toutes les personnes
qui ont contribué, de près ou de loin, à la réalisation de ce projet de
fin d'études.

\medskip
Nos remerciements les plus sincères s'adressent en premier lieu à notre
encadreur académique, \textbf{[Titre. Prénom~\textsc{Nom}]}, pour sa
disponibilité permanente, ses orientations rigoureuses, ses précieux
conseils et la confiance qu'il nous a accordée tout au long de
l'encadrement de ce travail.

\medskip
Nous remercions également notre encadreur en entreprise,
\textbf{[Titre. Prénom~\textsc{Nom}]}, pour son accueil, le partage de
son expertise en cybersécurité et en architecture logicielle, ainsi que
les recommandations déterminantes qu'il nous a prodiguées dans la conduite
de ce projet.

\medskip
Nos remerciements s'étendent à l'ensemble du corps enseignant et
administratif du~\textbf{[Nom de l'établissement]}, dont la qualité de
l'enseignement dispensé nous a fourni les fondements théoriques et
pratiques indispensables à la réalisation de ce travail.

\medskip
Nous remercions enfin tous nos collègues et amis pour leur soutien moral,
leurs suggestions constructives et leur bienveillance constante.

% =============================================================
% RÉSUMÉ (FRANÇAIS)
% =============================================================
\clearpage
\chapter*{Résumé}
\thispagestyle{plain}
\addcontentsline{toc}{chapter}{Résumé}
\markboth{Résumé}{Résumé}

Dans le cadre de notre Projet de Fin d'Études, nous avons conçu et
développé \textbf{AegisWeb}, une plateforme de gouvernance et de confiance
pour les agents d'intelligence artificielle agissant sur le web. Ce projet
répond à un défi émergent des entreprises modernes~: comment autoriser des
agents IA autonomes à naviguer sur des portails SaaS, télécharger des
factures, vérifier des renouvellements et effectuer des changements de
plan sans exposer de mots de passe bruts ni permettre d'actions
destructrices silencieuses.

\medskip
Notre solution repose sur une architecture monorepo TypeScript à quatre
composants applicatifs~: une API de contrôle \textbf{NestJS~11}, un worker
\textbf{BullMQ} avec runtime navigateur \textbf{Playwright}, un portail
SaaS simulé (\textit{vendor sandbox}), et un tableau de bord opérationnel
\textbf{Next.js~16}. La persistance est assurée par \textbf{PostgreSQL~17}
via \textbf{Prisma~7}, les files d'attente par \textbf{Redis~8}, et le
stockage de preuves par \textbf{MinIO} (S3-compatible). Le modèle de
sécurité intègre un coffre-fort AES-256-GCM, un moteur de politiques pur,
un RBAC à cinq rôles, des approbations humaines, une chaîne d'audit
SHA-256 et des reçus de confiance (\textit{receipts}).

\medskip
Les résultats obtenus attestent d'un MVP local complet et fonctionnel,
couvrant l'identité des agents, la gestion des fournisseurs SaaS, les
credentials chiffrés, les politiques de risque, trois workflows de
procurement (téléchargement de facture, vérification de renouvellement,
demande de downgrade avec approbation), ainsi qu'un dashboard de preuves
opérationnel. La validation repose sur 33 fichiers de tests Vitest
(195 tests passés) et des scripts E2E.

\medskip\noindent
\textbf{Mots-clés :}~\textit{agents IA, gouvernance, cybersécurité,
coffre-fort de secrets, moteur de politiques, approbation humaine,
Playwright, NestJS, Next.js, audit trail, procurement SaaS.}

% =============================================================
% ABSTRACT (ANGLAIS)
% =============================================================
\clearpage
\chapter*{Abstract}
\thispagestyle{plain}
\addcontentsline{toc}{chapter}{Abstract}
\markboth{Abstract}{Abstract}

As part of our Final Year Project, we designed and developed
\textbf{AegisWeb}, a governance and trust platform for AI agents acting
on the web. This project addresses a critical challenge for modern
organizations: how to let autonomous AI agents browse SaaS portals,
download invoices, check renewals, and perform plan changes without
exposing raw credentials or allowing silent destructive actions.

\medskip
Our solution is built upon a TypeScript monorepo with four application
components: a \textbf{NestJS~11} control-plane API, a \textbf{BullMQ}
worker with \textbf{Playwright} browser runtime, a simulated SaaS vendor
sandbox, and a \textbf{Next.js~16} operational dashboard. Persistence
relies on \textbf{PostgreSQL~17} via \textbf{Prisma~7}, queuing on
\textbf{Redis~8}, and evidence storage on \textbf{MinIO} (S3-compatible).
The security model includes an AES-256-GCM vault, a pure policy engine,
five-role RBAC, human approvals, SHA-256 audit hash chains, and trust
receipts.

\medskip
The results demonstrate a complete and functional local MVP covering agent
identity, SaaS vendor management, encrypted credentials, risk policies,
three procurement workflows (invoice download, renewal check, downgrade
request with approval), and an evidence-first operational dashboard.
Validation relies on 33 Vitest test files (195 passing tests) and E2E
scripts.

\medskip\noindent
\textbf{Keywords:}~\textit{AI agents, governance, cybersecurity, secret
vault, policy engine, human approval, Playwright, NestJS, Next.js,
audit trail, SaaS procurement.}

% =============================================================
% TABLE DES MATIÈRES
% =============================================================
\clearpage
\tableofcontents
\thispagestyle{plain}

% =============================================================
% LISTE DES FIGURES
% =============================================================
\clearpage
\listoffigures
\addcontentsline{toc}{chapter}{Liste des figures}

% =============================================================
% LISTE DES TABLEAUX
% =============================================================
\clearpage
\listoftables
\addcontentsline{toc}{chapter}{Liste des tableaux}

% =============================================================
% LISTE DES ABRÉVIATIONS
% =============================================================
\clearpage
\chapter*{Liste des abréviations}
\thispagestyle{plain}
\addcontentsline{toc}{chapter}{Liste des abréviations}
\markboth{Liste des abréviations}{Liste des abréviations}

\vspace{4pt}
\begin{tabular}{@{}p{2.8cm}p{11.5cm}@{}}
\toprule
\textbf{Abréviation} & \textbf{Signification} \\
\midrule
AES-GCM    & Advanced Encryption Standard — Galois/Counter Mode \\
AI / IA    & Artificial Intelligence / Intelligence Artificielle \\
API        & Application Programming Interface \\
Argon2id   & Algorithme de hachage de mots de passe \\
BFF        & Backend For Frontend \\
BullMQ     & File de messages Redis pour Node.js \\
CORS       & Cross-Origin Resource Sharing \\
CRUD       & Create, Read, Update, Delete \\
CSP        & Content Security Policy \\
DI         & Dependency Injection \\
DTO        & Data Transfer Object \\
GCM        & Galois/Counter Mode (mode de chiffrement) \\
HMAC       & Hash-based Message Authentication Code \\
HTTP/S     & Hypertext Transfer Protocol (Secure) \\
IDE        & Integrated Development Environment \\
JWT        & JSON Web Token \\
MVP        & Minimum Viable Product \\
ORM        & Object-Relational Mapping \\
PFE        & Projet de Fin d'Études \\
RBAC       & Role-Based Access Control \\
REST       & Representational State Transfer \\
RPA        & Robotic Process Automation \\
S3         & Simple Storage Service (stockage objet) \\
SaaS       & Software as a Service \\
SDK        & Software Development Kit \\
SMTP       & Simple Mail Transfer Protocol \\
SQL        & Structured Query Language \\
TLS        & Transport Layer Security \\
TTL        & Time To Live \\
UML        & Unified Modeling Language \\
URL        & Uniform Resource Locator \\
UUID       & Universally Unique Identifier \\
\bottomrule
\end{tabular}

% =============================================================
%      CORPS DU RAPPORT
% =============================================================
\cleardoublepage
\pagenumbering{arabic}
\setcounter{page}{1}

% =============================================================
% INTRODUCTION GÉNÉRALE
% =============================================================
\clearpage
\chapter*{Introduction Générale}
\addcontentsline{toc}{chapter}{Introduction Générale}
\markboth{Introduction Générale}{Introduction Générale}

\begin{chapintro}
L'émergence des agents d'intelligence artificielle capables d'agir de
manière autonome sur le web transforme profondément les processus
opérationnels des entreprises. Ces agents peuvent désormais naviguer sur
des portails SaaS, remplir des formulaires, télécharger des factures et
modifier des paramètres de facturation. Pourtant, le goulot d'étranglement
n'est plus seulement l'intelligence~: c'est la \textbf{permission}.
\end{chapintro}

Les entreprises ne peuvent pas confier à des agents autonomes des mots de
passe bruts, un accès administrateur illimité, une autorité de paiement
non contrôlée ou un contrôle navigateur sans garde-fous. Les solutions
RPA traditionnelles et les intégrations API classiques ne couvrent pas le
besoin de gouvernance fine, d'approbation humaine contextuelle et de
preuves vérifiables pour chaque action sensible.

\medskip
\noindent\textbf{Problématique :} Comment concevoir et développer une
plateforme de gouvernance pour agents IA web, offrant des identités
dédiées, un coffre-fort de credentials, un moteur de politiques, des
approbations humaines, une exécution navigateur contrôlée et un audit
traçable, tout en répondant aux exigences de sécurité, de performance et
de maintenabilité d'un système distribué moderne ?

\medskip
\noindent\textbf{Organisation du rapport :} Ce rapport s'articule autour
de quatre chapitres complémentaires :

\begin{itemize}[leftmargin=2cm, itemsep=3pt]
  \item \textbf{Chapitre~1 — Cadre du projet :} présentation du contexte,
    analyse de l'existant, description de la solution AegisWeb et
    méthodologie de développement incrémental.
  \item \textbf{Chapitre~2 — Analyse et spécification des besoins :}
    identification des acteurs, besoins fonctionnels et non
    fonctionnels, cas d'utilisation et backlog produit.
  \item \textbf{Chapitre~3 — Conception du système :} architecture
    monorepo, modélisations dynamique et statique, dictionnaire de
    données.
  \item \textbf{Chapitre~4 — Réalisation :} environnement, technologies,
    livrables par phase, interfaces et résultats des tests.
\end{itemize}

% =============================================================
% CHAPITRE 1 — CADRE DU PROJET
% =============================================================
\chapter{Cadre du projet}

\begin{chapintro}
Ce premier chapitre situe notre projet dans son contexte professionnel
et technique. Nous y présentons le domaine de la gouvernance des agents
IA, analysons l'existant pour en identifier les lacunes, décrivons la
solution AegisWeb, puis exposons la méthodologie et le planning retenus.
\end{chapintro}

\section{Présentation du contexte et de l'organisme d'accueil}

\subsection{Contexte du projet}

\textbf{AegisWeb} est une plateforme de \textbf{gouvernance et de confiance
pour agents IA agissant sur le web}. Elle fournit une couche de contrôle
entre les agents autonomes et les systèmes métier (portails SaaS,
formulaires, facturation, etc.). Son principe fondateur est
\textit{«~Permission before autonomy~»}~: la permission est le goulot
d'étranglement, pas l'intelligence.

Le premier cas d'usage ciblé est le \textbf{procurement SaaS}~: permettre
à des agents IA de gérer en toute sécurité les factures, les
renouvellements et les changements de plan soumis à approbation humaine,
sans exposer de credentials bruts.

\begin{table}[H]
\centering
\caption{Fiche d'identité du projet AegisWeb}
\label{tab:fiche_projet}
\begin{tabularx}{\textwidth}{@{}lX@{}}
\toprule
\textbf{Critère}         & \textbf{Description} \\
\midrule
Dénomination             & AegisWeb (nom interne historique~: AgentPass) \\
\rowcolor{rowgray}
Secteur d'activité       & Cybersécurité applicative, gouvernance IA \\
Nature                   & Plateforme MVP locale de gouvernance d'agents web \\
\rowcolor{rowgray}
Produit phare            & Couche permission, credential, approbation, runtime
                           navigateur et audit pour agents IA \\
Marché cible             & Entreprises utilisant des agents IA pour
                           l'automatisation web contrôlée \\
\rowcolor{rowgray}
Cas d'usage MVP          & Procurement SaaS (factures, renouvellements,
                           downgrades avec approbation) \\
Langues de l'interface   & Anglais (dashboard), documentation bilingue \\
\bottomrule
\end{tabularx}
\end{table}

\subsection{Cadre du projet de fin d'études}

Notre projet de fin d'études consiste en la \textbf{conception et le
développement complets de la plateforme AegisWeb}, couvrant~:
\begin{itemize}[leftmargin=1.8cm, itemsep=3pt]
  \item le \textbf{backend} (API REST NestJS, base de données, files
    BullMQ, endpoints internes worker) ;
  \item le \textbf{worker} (exécution de workflows, Playwright, connecteur
    vendor) ;
  \item le \textbf{frontend} (dashboard Next.js, authentification, couche
    données React Query) ;
  \item les \textbf{bibliothèques partagées} (domain, vault, policy-engine,
    browser-runtime, audit) ;
  \item l'\textbf{infrastructure} locale (Docker Compose, PostgreSQL,
    Redis, MinIO, Mailpit).
\end{itemize}

\section{Étude de l'existant}

\subsection{Présentation des solutions existantes}

\begin{table}[H]
\centering
\caption{Analyse comparative des approches existantes}
\label{tab:existant}
\begin{tabularx}{\textwidth}{@{}lXXC{2.4cm}@{}}
\toprule
\textbf{Solution} & \textbf{Points forts} & \textbf{Limitations}
  & \textbf{Gouvernance IA} \\
\midrule
\rowcolor{rowgray}
RPA classique (UiPath, etc.)
  & Automatisation UI mature, intégrations entreprise
  & Pas d'identité agent, pas de politiques fines, audit limité
  & Faible \\
Intégrations API directes
  & Rapides, déterministes, testables
  & Couverture partielle des portails, pas d'actions UI
  & Moyenne \\
\rowcolor{rowgray}
Gestionnaires de secrets (Vault)
  & Chiffrement robuste, rotation, audit
  & Pas de runtime navigateur, pas d'approbation workflow
  & Partielle \\
Agents IA autonomes (bruts)
  & Flexibilité maximale, intelligence adaptative
  & Risque credential, actions silencieuses, pas de preuves
  & Inexistante \\
\rowcolor{rowgray}
\textbf{AegisWeb (proposé)}
  & Identité agent, vault, politiques, approbation, preuves, audit
  & MVP local, pas encore SaaS production multi-tenant
  & Élevée \\
\bottomrule
\end{tabularx}
\end{table}

\subsection{Critique de l'existant}

L'analyse révèle plusieurs lacunes fondamentales :
\begin{itemize}[leftmargin=1.8cm, itemsep=4pt]
  \item \textbf{Absence d'identité non-humaine dédiée} pour les agents,
    avec statut, purpose et grants explicites.
  \item \textbf{Exposition des credentials} aux agents ou aux prompts LLM,
    sans chiffrement au repos ni déchiffrement scopé par exécution.
  \item \textbf{Absence de moteur de politiques} évaluant domaines,
    actions interdites, seuils financiers et niveaux de risque.
  \item \textbf{Pas de pause pour approbation humaine} avant les actions
    destructrices (downgrade, annulation).
  \item \textbf{Manque de preuves opérationnelles}~: screenshots, factures,
    reçus de confiance et chaîne d'audit vérifiable.
  \item \textbf{Runtime navigateur non contrôlé}~: accès réseau privé,
    popups hors domaine, téléchargements non tracés.
\end{itemize}

\section{Solution proposée}

\subsection{Description d'AegisWeb}

\textbf{AegisWeb} est la couche de permission, credential, approbation,
runtime navigateur et audit pour les agents IA agissant sur le web. Elle
donne aux agents une \textbf{autorité contrôlée} via~:

\begin{itemize}[leftmargin=1.8cm, itemsep=4pt]
  \item \textbf{Identité agent}~: agent non-humain avec identifier unique,
    statut (active, paused, revoked) et purpose documenté.
  \item \textbf{Coffre-fort de credentials}~: chiffrement AES-256-GCM au
    repos, octroi par agent via \texttt{CredentialAgentGrant}.
  \item \textbf{Moteur de politiques}~: allowlist de domaines, actions
    interdites, seuils de dépense, règles d'approbation.
  \item \textbf{Runtime navigateur contrôlé}~: Playwright avec allowlist,
    blocage réseau privé, masquage des champs sensibles.
  \item \textbf{Workflows queue-backed}~: trois templates SaaS procurement.
  \item \textbf{Approbation humaine}~: pause du run, notification email,
    reprise conditionnelle.
  \item \textbf{Audit trail}~: événements append-only avec chaîne de hash
    SHA-256 (\texttt{prevHash} + \texttt{eventHash}).
  \item \textbf{Receipts}~: artefacts finaux de confiance résumant
    timeline, décisions et preuves.
  \item \textbf{Dashboard opérationnel}~: salle de preuves pour agents,
    vendors, policies, runs, approvals, audit.
\end{itemize}

\begin{figure}[H]
\centering
\includegraphics[width=0.55\textwidth]{assets/aegis_agent_identity_passport.png}
\caption{Concept d'identité agent AegisWeb — passeport d'autorité contrôlée}
\label{fig:agent_passport}
\end{figure}

\subsection{Flux produit phare — Downgrade de plan SaaS}

\begin{enumerate}[leftmargin=1.8cm, itemsep=3pt]
  \item L'owner se connecte au dashboard.
  \item Il crée ou sélectionne agent, vendor, credential, policy, workflow.
  \item Il lance «~Acme Downgrade Request~».
  \item L'API crée un \texttt{WorkflowRun} et enqueue un job BullMQ.
  \item Le worker ouvre un navigateur contrôlé (Playwright).
  \item Le worker se connecte au sandbox via credential déchiffré (API
    interne vault, scopé au run).
  \item Le worker lit les données de renouvellement et de plan.
  \item Le moteur de politique classe l'action comme risquée
    $\rightarrow$ \textbf{pause pour approbation}.
  \item L'approbateur examine les preuves (screenshots) et approuve ou
    rejette.
  \item Le worker reprend uniquement après approbation.
  \item L'action est soumise $\rightarrow$ \textbf{receipt} + \textbf{audit
    trail} générés.
  \item Notification email via Mailpit (en local).
\end{enumerate}

\subsection{Objectifs du projet}

\paragraph{Objectifs fonctionnels :}
\begin{itemize}[leftmargin=1.8cm, itemsep=3pt]
  \item Développer une API REST NestJS avec 20+ modules (auth, RBAC,
    agents, vendors, policies, credentials, workflows, runs, approvals,
    receipts, audit, files, queue, internal worker).
  \item Implémenter un schéma Prisma multi-tenant (17 tables) avec seed
    de démonstration riche.
  \item Concevoir un moteur de politiques pur TypeScript sans dépendance
    BDD.
  \item Réaliser un coffre-fort AES-256-GCM avec redaction récursive des
    secrets.
  \item Développer un worker BullMQ + Playwright avec trois workflows
    SaaS.
  \item Mettre en place un dashboard Next.js 16 evidence-first.
  \item Fournir un vendor sandbox pour démos et tests déterministes.
\end{itemize}

\paragraph{Objectifs non fonctionnels :}
\begin{itemize}[leftmargin=1.8cm, itemsep=3pt]
  \item Isolation multi-tenant stricte par \texttt{organizationId}.
  \item Argon2id, JWT courte durée, refresh HttpOnly atomique.
  \item Rate limiting Redis en production.
  \item CSP nonce-based sans \texttt{unsafe-eval} sur le frontend.
  \item Tests d'intégration Vitest couvrant les 30 phases backend.
  \item Déploiement Docker Compose reproductible.
\end{itemize}

\subsection{Périmètre du projet}

Le périmètre couvre~:
\begin{enumerate}[leftmargin=1.8cm, itemsep=3pt]
  \item Le monorepo \texttt{agentpass} (pnpm workspaces).
  \item Les applications \texttt{api}, \texttt{worker}, \texttt{web},
    \texttt{vendor-sandbox}.
  \item Les bibliothèques \texttt{domain}, \texttt{database}, \texttt{vault},
    \texttt{policy-engine}, \texttt{browser-runtime}, \texttt{audit},
    \texttt{testing}.
  \item L'infrastructure \texttt{infra/docker-compose.yml}.
  \item 33 fichiers de tests Vitest + scripts smoke/E2E.
\end{enumerate}

\section{Méthodologie de gestion de projet}

\subsection{Approche incrémentale par phases}

Contrairement à une approche SCRUM classique par sprints, AegisWeb a été
développé selon une \textbf{méthodologie incrémentale en 30 phases}
(documentées dans \texttt{docs/BACKEND\_IMPLEMENTATION\_PLAN.md}), chaque
phase livrant un module testable et validé avant passage à la suivante.

\begin{table}[H]
\centering
\caption{Principes de la méthodologie par phases}
\label{tab:methodo}
\begin{tabularx}{\textwidth}{@{}lX@{}}
\toprule
\textbf{Élément}        & \textbf{Application dans le projet AegisWeb} \\
\midrule
\rowcolor{rowgray}
Phase                   & Incrément fonctionnel avec tests dédiés
                          (\texttt{phaseN.spec.ts}) \\
Definition of Done      & Tests passants, module intégré, notes de
                          complétion documentées \\
\rowcolor{rowgray}
Acceptance MVP          & \texttt{mvp-backend-acceptance.spec.ts}
                          (31 fichiers, 195 tests) \\
Traçabilité             & \texttt{docs/PHASE\_*\_COMPLETION\_NOTES.md}
                          (27 fichiers) \\
\rowcolor{rowgray}
Sécurité continue       & \texttt{security-coverage.spec.ts} +
                          \texttt{docs/SECURITY\_AUDIT.md} \\
\bottomrule
\end{tabularx}
\end{table}

\subsection{Planification générale — Diagramme de Gantt}

\begin{figure}[H]
\centering
\caption{Diagramme de Gantt — Planning prévisionnel du projet (24 semaines)}
\label{fig:gantt}
\resizebox{\textwidth}{!}{%
\begin{ganttchart}[
  hgrid,
  vgrid          = {*{1}{dotted}},
  bar/.style     = {fill=aegisblue!70, rounded corners=1pt},
  bar label font = \small,
  title label font = \small\bfseries,
  milestone label font = \small\itshape,
  milestone/.style  = {fill=aegisaccent},
  x unit         = 0.45cm,
  y unit title   = 0.62cm,
  y unit chart   = 0.68cm,
  title height   = 1,
  bar height     = 0.55,
]{1}{24}
  \gantttitle{Planning prévisionnel — AegisWeb (24 semaines)}{24} \\
  \gantttitlelist{1,...,24}{1} \\
  \ganttbar{Étude \& Analyse besoins}{1}{2} \\
  \ganttbar{Phases 0--3 : Fondation API}{3}{5} \\
  \ganttbar{Phases 4--8 : Auth, RBAC, Org, Users}{6}{8} \\
  \ganttbar{Phases 9--13 : Agents, Vendors, Vault}{9}{11} \\
  \ganttbar{Phases 14--18 : Workflows, Queue}{12}{14} \\
  \ganttbar{Phases 19--23 : Worker, Browser, Sandbox}{15}{17} \\
  \ganttbar{Phases 24--26 : Workflows SaaS}{18}{19} \\
  \ganttbar{Phases 27--29 : Receipts, Notif, OpenAPI}{20}{21} \\
  \ganttbar{Dashboard Next.js (frontend)}{10}{22} \\
  \ganttbar{Tests acceptance \& sécurité}{21}{23} \\
  \ganttbar{Rédaction rapport PFE}{22}{24} \\
  \ganttmilestone{Livraison finale}{24}
\end{ganttchart}%
}
\end{figure}

\conclusionchap{Ce chapitre nous a permis de situer AegisWeb dans le
contexte de la gouvernance des agents IA web, d'identifier les
insuffisances des solutions existantes et de définir les contours de
notre plateforme. La méthodologie incrémentale en 30 phases et le planning
sur 24 semaines constituent le cadre organisationnel de notre démarche.
Le chapitre suivant sera consacré à l'analyse détaillée des besoins.}

% =============================================================
% CHAPITRE 2 — ANALYSE ET SPÉCIFICATION DES BESOINS
% =============================================================
\chapter{Analyse et spécification des besoins}

\begin{chapintro}
Ce chapitre est dédié à l'analyse des besoins de la plateforme AegisWeb.
Nous y identifions les acteurs, détaillons les besoins fonctionnels et
non fonctionnels, présentons les diagrammes de cas d'utilisation et
définissons le backlog produit initial.
\end{chapintro}

\section{Identification des acteurs}

\begin{table}[H]
\centering
\caption{Tableau des acteurs du système AegisWeb}
\label{tab:acteurs}
\begin{tabularx}{\textwidth}{@{}llX@{}}
\toprule
\textbf{Acteur}             & \textbf{Type}   & \textbf{Rôle et description} \\
\midrule
\rowcolor{rowgray}
Owner / Admin               & Primaire
  & Propriétaire de l'organisation. Gère agents, vendors, policies,
    credentials, workflows. Lance et annule des runs. \\
Approver                    & Primaire
  & Examine les demandes d'approbation, consulte les preuves (screenshots,
    audit) et approuve ou rejette les actions risquées. \\
\rowcolor{rowgray}
Developer                   & Primaire
  & Développeur technique. Consulte les ressources et peut lancer/annuler
    des workflows (\texttt{workflow:run}, \texttt{workflow:cancel}). \\
Auditor                     & Primaire
  & Auditeur de conformité. Accès lecture sur agents, runs, receipts,
    audit trail et fichiers de preuve. \\
\rowcolor{rowgray}
Agent IA                    & Secondaire
  & Identité non-humaine exécutant des workflows via le worker, soumise
    aux politiques et grants de credentials. \\
Worker                      & Secondaire
  & Processus BullMQ consommant les jobs, exécutant Playwright et
    appelant l'API interne avec token HMAC scopé. \\
\rowcolor{rowgray}
Vendor Sandbox              & Secondaire
  & Portail SaaS simulé (facturation, factures, downgrade) pour démos
    et tests sans credentials réels. \\
Système (Mailpit/S3)        & Secondaire
  & Services d'infrastructure~: stockage objet MinIO, notifications SMTP
    Mailpit, Redis BullMQ. \\
\bottomrule
\end{tabularx}
\end{table}

\section{Besoins fonctionnels}

\begin{longtable}{@{}p{1.2cm}p{2.2cm}p{10.2cm}@{}}
\caption{Besoins fonctionnels de la plateforme AegisWeb}
\label{tab:bf} \\
\toprule
\textbf{ID} & \textbf{Module} & \textbf{Description} \\
\midrule
\endfirsthead
\toprule
\textbf{ID} & \textbf{Module} & \textbf{Description} \\
\midrule
\endhead
\bottomrule
\endfoot
\bottomrule
\endlastfoot
BF01 & Auth.        & Inscription organisation, login Argon2id, JWT access
                      token, refresh token HttpOnly (\texttt{agentpass\_refresh\_token}),
                      refresh atomique, logout. \\
\rowcolor{rowgray}
BF02 & RBAC         & Cinq rôles (owner, admin, approver, auditor, developer)
                      avec 28 permissions granulaires et isolation
                      \texttt{organizationId}. \\
BF03 & Agents       & CRUD agents non-humains (identifier, purpose, status,
                      grants credentials). \\
\rowcolor{rowgray}
BF04 & Vendors      & Gestion fournisseurs SaaS (website, catégorie, coût
                      mensuel estimé). \\
BF05 & Policies     & Politiques JSON typées (allowlist, actions, spending,
                      approval rules, bundles). \\
\rowcolor{rowgray}
BF06 & Vault        & Credentials chiffrés AES-256-GCM, octroi/révocation
                      par agent, déchiffrement scopé run. \\
BF07 & Workflows    & Définitions réutilisables (3 templates)~:
                      invoice download, renewal check, plan downgrade. \\
\rowcolor{rowgray}
BF08 & Runs         & Exécutions workflow avec états (queued, running,
                      waiting\_for\_approval, completed, failed, canceled, denied). \\
BF09 & Policy Eng.  & Évaluation pure~: agent actif, domaine, allowlist,
                      actions interdites, seuils, approbation, risque. \\
\rowcolor{rowgray}
BF10 & Browser RT   & Playwright contrôlé~: allowlist domaines, blocage
                      réseau privé, screenshots masqués, downloads SHA-256. \\
BF11 & Approvals    & Pause run, création \texttt{ApprovalRequest},
                      notification email, approve/reject, reprise queue. \\
\rowcolor{rowgray}
BF12 & Audit        & Événements typés append-only, chaîne hash SHA-256,
                      acteurs (user, agent, worker, system). \\
BF13 & Files        & Métadonnées fichiers S3 (screenshots, factures PDF,
                      traces) avec isolation org. \\
\rowcolor{rowgray}
BF14 & Receipts     & Artefacts finaux de confiance (timeline, décisions,
                      approbation, preuves). \\
BF15 & Dashboard    & Interface Next.js pour toutes les ressources +
                      panneaux evidence (timeline, screenshots, audit drawer). \\
\rowcolor{rowgray}
BF16 & Sandbox      & Portail SaaS factice (login, billing, invoices,
                      downgrade, cancel, admin users). \\
BF17 & Queue        & BullMQ~: \texttt{workflow-runs}, \texttt{workflow-resume},
                      \texttt{workflow-maintenance} avec jobs idempotents. \\
\rowcolor{rowgray}
BF18 & OpenAPI      & Documentation Swagger activable via
                      \texttt{ENABLE\_OPENAPI=true}. \\
\end{longtable}

\section{Besoins non fonctionnels}

\begin{table}[H]
\centering
\caption{Besoins non fonctionnels de la plateforme AegisWeb}
\label{tab:bnf}
\begin{tabularx}{\textwidth}{@{}p{2.5cm}X@{}}
\toprule
\textbf{Critère}     & \textbf{Exigence} \\
\midrule
\rowcolor{rowgray}
Sécurité             & Argon2id, JWT courte durée, refresh HttpOnly atomique,
                       AES-256-GCM vault, tokens worker HMAC, CSP nonce,
                       rate limiting Redis, redaction secrets dans audit/receipts. \\
Isolation tenant     & Toutes les requêtes API scopées par
                       \texttt{organizationId}. Tests régression cross-org. \\
\rowcolor{rowgray}
Traçabilité          & Audit append-only avec \texttt{prevHash} et
                       \texttt{eventHash}. Receipts immuables post-exécution. \\
Performance          & Worker concurrency = 1 par défaut. Retry BullMQ
                       3 tentatives, backoff exponentiel. \\
\rowcolor{rowgray}
Maintenabilité       & Monorepo TypeScript, libs partagées, Prisma migrations,
                       33 specs Vitest, ESLint 9. \\
Disponibilité        & Endpoints \texttt{/health} et \texttt{/health/ready}
                       avec vérification Postgres/Redis. \\
\rowcolor{rowgray}
Déployabilité        & Docker multi-stage (api, worker, web, vendor-sandbox),
                       profils infra et app, guide Render. \\
UX Dashboard         & Design evidence-first, dense, monochrome, orienté
                       confiance. Aucun secret en clair dans l'UI. \\
\rowcolor{rowgray}
Pilot readiness      & Score audit sécurité documenté~: 3,8/5 global
                       (\texttt{docs/SECURITY\_AUDIT.md}). \\
\bottomrule
\end{tabularx}
\end{table}

\section{Diagrammes de cas d'utilisation}

\subsection{Diagramme de cas d'utilisation global}

\begin{figure}[H]
\centering
\diagramespace[7.5cm]{Diagramme CU global — acteurs Owner, Approver,
  Developer, Auditor, Agent IA, Worker — CU~: S'authentifier, Gérer
  agents/vendors/policies/credentials, Lancer workflow, Approuver action,
  Consulter audit/receipts, Exécuter run navigateur}
\caption{Diagramme de cas d'utilisation global du système AegisWeb}
\label{fig:cu_global}
\end{figure}

\subsection{Module Authentification et autorisation}

\begin{figure}[H]
\centering
\diagramespace[5.5cm]{Diagramme CU — Auth (S'inscrire, Se connecter,
  Rafraîchir token, Se déconnecter, Consulter profil) + RBAC par rôle}
\caption{Diagramme CU — Module Authentification et RBAC}
\label{fig:cu_auth}
\end{figure}

\subsection{Module Workflow et approbation}

\begin{figure}[H]
\centering
\diagramespace[6cm]{Diagramme CU — Workflow (Créer workflow, Lancer run,
  [extend] Demander approbation, Approuver/Rejeter, Reprendre run,
  Générer receipt) — acteurs Owner, Approver, Worker}
\caption{Diagramme CU — Module Workflow et approbation}
\label{fig:cu_workflow}
\end{figure}

\section{Descriptions textuelles des cas d'utilisation}

\subsection{CU — Lancer un workflow de downgrade avec approbation}

\begin{tcolorbox}[aegiscubox,
  title=Cas d'utilisation CU-WF-01 : Lancer un downgrade SaaS avec approbation]
\begin{tabularx}{\textwidth}{@{}p{3.2cm}X@{}}
\textbf{Acteur principal}    & Owner ou Developer \\
\textbf{Acteurs secondaires} & Worker, Policy Engine, Approver, Mailpit \\
\textbf{Pré-condition}       & Agent actif, credential accordé, policy active,
                               workflow \texttt{plan\_downgrade\_request} configuré \\
\textbf{Post-condition}      & Run \texttt{completed} ou \texttt{denied}, receipt
                               et audit générés \\
\textbf{Déclencheur}         & Clic «~Run workflow~» sur le dashboard \\
\end{tabularx}

\smallskip
\textbf{Scénario nominal :}
\begin{enumerate}[leftmargin=1.2cm, itemsep=2pt]
  \item L'utilisateur authentifié lance le workflow depuis le dashboard.
  \item L'API crée un \texttt{WorkflowRun} (statut \texttt{queued}) et
        enqueue un job \texttt{start-\{runId\}}.
  \item Le worker consomme le job, passe le run en \texttt{running}.
  \item Le worker ouvre un contexte Playwright isolé (allowlist vendor).
  \item Le worker demande le déchiffrement du credential via
        \texttt{POST /internal/vault/credentials/:id/decrypt-for-run}.
  \item Le worker se connecte au vendor sandbox et lit les données billing.
  \item Chaque action est enregistrée comme \texttt{ActionAttempt} avec
        décision du policy engine.
  \item Pour \texttt{change\_plan}~: décision \texttt{require\_approval}
        $\rightarrow$ run \texttt{waiting\_for\_approval}.
  \item Email envoyé à l'approver via Mailpit.
  \item L'approver examine screenshots et approuve.
  \item Job \texttt{resume-\{runId\}-\{approvalId\}} enqueue.
  \item Le worker soumet le downgrade et génère receipt + audit.
\end{enumerate}

\smallskip
\textbf{Scénarios alternatifs :}
\begin{itemize}[leftmargin=1.2cm, itemsep=2pt]
  \item \textbf{7a.} Décision \texttt{deny} $\rightarrow$ run
        \texttt{denied}, pas de soumission.
  \item \textbf{10a.} Rejet approbation $\rightarrow$ run \texttt{denied}.
  \item \textbf{E1.} Credential révoqué $\rightarrow$ run \texttt{failed}.
  \item \textbf{E2.} Domaine hors allowlist $\rightarrow$ action bloquée
        par browser-runtime.
\end{itemize}
\end{tcolorbox}

\subsection{CU — Créer et accorder un credential}

\begin{tcolorbox}[aegiscubox,
  title=Cas d'utilisation CU-VAULT-01 : Créer et accorder un credential]
\begin{tabularx}{\textwidth}{@{}p{3.2cm}X@{}}
\textbf{Acteur principal}    & Owner ou Admin \\
\textbf{Acteurs secondaires} & Vault (AES-256-GCM), Audit \\
\textbf{Pré-condition}       & Vendor existant, permission
                               \texttt{credential:create} \\
\textbf{Post-condition}      & Credential chiffré créé, grant agent optionnel,
                               événement audit enregistré \\
\textbf{Déclencheur}         & Formulaire «~New credential~» sur le dashboard \\
\end{tabularx}

\smallskip
\textbf{Scénario nominal :}
\begin{enumerate}[leftmargin=1.2cm, itemsep=2pt]
  \item L'utilisateur saisit le type (username/password, API token, etc.)
        et le payload secret.
  \item L'API chiffre le payload via \texttt{encryptSecret} (vault lib).
  \item Persistance Prisma~: seul le ciphertext est stocké.
  \item L'utilisateur accorde l'accès à un agent
        (\texttt{CredentialAgentGrant}).
  \item Événements \texttt{credential\_created} et
        \texttt{credential\_granted\_to\_agent} dans l'audit trail.
  \item L'UI affiche les métadonnées sans jamais exposer le secret en clair.
\end{enumerate}
\end{tcolorbox}

\section{Backlog produit — Phases principales}

\begin{longtable}{@{}cp{5.8cm}p{2cm}cc@{}}
\caption{Backlog produit — Phases principales du développement}
\label{tab:backlog} \\
\toprule
\textbf{ID} & \textbf{Livrable}
  & \textbf{Priorité} & \textbf{Phase} & \textbf{Statut} \\
\midrule
\endfirsthead
\toprule
\textbf{ID} & \textbf{Livrable}
  & \textbf{Priorité} & \textbf{Phase} & \textbf{Statut} \\
\midrule
\endhead
\bottomrule
\endfoot
\bottomrule
\endlastfoot
P01 & Fondation API, worker boot, libs partagées     & Must Have & 0  & \coche \\
\rowcolor{rowgray}
P02 & Schéma Prisma + seed démo Northstar Labs        & Must Have & 1  & \coche \\
P03 & Auth Argon2id + JWT + refresh HttpOnly          & Must Have & 4  & \coche \\
\rowcolor{rowgray}
P04 & RBAC 5 rôles + isolation organisation           & Must Have & 5  & \coche \\
P05 & Audit append-only + hash chain                  & Must Have & 6  & \coche \\
\rowcolor{rowgray}
P06 & Files S3/MinIO + métadonnées                    & Must Have & 7  & \coche \\
P07 & Agents, Vendors, Policies CRUD                  & Must Have & 9--12 & \coche \\
\rowcolor{rowgray}
P08 & Vault AES-GCM + CredentialsModule               & Must Have & 13 & \coche \\
P09 & Workflows + WorkflowRuns + Queue BullMQ         & Must Have & 14--18 & \coche \\
\rowcolor{rowgray}
P10 & Worker + Browser Runtime Playwright               & Must Have & 19--20 & \coche \\
P11 & Vendor Sandbox + Connector                        & Must Have & 21--22 & \coche \\
\rowcolor{rowgray}
P12 & 3 workflows SaaS (invoice, renewal, downgrade)    & Must Have & 24--26 & \coche \\
P13 & Approvals + Notifications email                   & Must Have & 17, 28 & \coche \\
\rowcolor{rowgray}
P14 & Receipts + preuves finales                        & Must Have & 27 & \coche \\
P15 & Dashboard Next.js complet                         & Must Have & —  & \coche \\
\rowcolor{rowgray}
P16 & OpenAPI / Swagger                                 & Should Have & 29 & \coche \\
P17 & Tests acceptance MVP (195 tests)                  & Must Have & —  & \coche \\
\end{longtable}

\conclusionchap{Ce chapitre a permis d'identifier les 8 catégories d'acteurs
du système AegisWeb et de formaliser 18 besoins fonctionnels et 9 critères
non fonctionnels. Le backlog de 17 livrables majeurs, aligné sur les 30
phases de développement, constitue la feuille de route de notre
implémentation. Le chapitre suivant sera consacré à la conception
architecturale et à la modélisation détaillée du système.}

% =============================================================
% CHAPITRE 3 — CONCEPTION DU SYSTÈME
% =============================================================
\chapter{Conception du système}

\begin{chapintro}
Ce chapitre décrit la conception architecturale et la modélisation du
système AegisWeb. Nous y présentons l'architecture monorepo, les
modélisations dynamiques et statiques, ainsi que le dictionnaire de
données Prisma.
\end{chapintro}

\section{Architecture générale du système}

\subsection{Architecture logicielle — Monorepo}

AegisWeb repose sur une \textbf{architecture monorepo TypeScript} (pnpm
workspaces) avec séparation claire entre applications et bibliothèques
partagées.

\begin{figure}[H]
\centering
\diagramespace[6.5cm]{Schéma d'architecture AegisWeb — Utilisateur/Approver
  $\rightarrow$ Next.js Web (3000) $\rightarrow$ NestJS API (3001)
  $\rightarrow$ PostgreSQL + Redis/BullMQ + MinIO/S3 — Worker
  $\rightarrow$ Playwright $\rightarrow$ Vendor Sandbox (4202)}
\caption{Architecture logicielle globale d'AegisWeb}
\label{fig:archi_global}
\end{figure}

\begin{table}[H]
\centering
\caption{Composants applicatifs et bibliothèques partagées}
\label{tab:composants}
\begin{tabularx}{\textwidth}{@{}p{2.8cm}Xp{4.2cm}@{}}
\toprule
\textbf{Composant}   & \textbf{Rôle}                              & \textbf{Technologies} \\
\midrule
\rowcolor{rowgray}
\texttt{apps/web}    & Dashboard + marketing, auth session        & Next.js 16, React 19, Tailwind 4 \\
\texttt{apps/api}    & Control plane REST, queue, internal API    & NestJS 11, Prisma 7, BullMQ \\
\rowcolor{rowgray}
\texttt{apps/worker} & Exécution workflows, Playwright            & NestJS context, BullMQ \\
\texttt{apps/vendor-sandbox} & Portail SaaS simulé                 & NestJS minimal \\
\rowcolor{rowgray}
\texttt{libs/domain} & Enums, permissions, types queue/workflow   & TypeScript pur \\
\texttt{libs/vault}  & Chiffrement AES-256-GCM, redaction         & Node.js crypto \\
\rowcolor{rowgray}
\texttt{libs/policy-engine} & Évaluation politiques/risk          & TypeScript pur \\
\texttt{libs/browser-runtime} & Wrapper Playwright contrôlé       & Playwright Chromium \\
\bottomrule
\end{tabularx}
\end{table}

\subsection{Modules NestJS de l'API}

L'API (\texttt{apps/api/src/app.module.ts}) importe 20+ modules~:

\begin{table}[H]
\centering
\caption{Modules principaux de l'API NestJS}
\label{tab:modules_api}
\small
\begin{tabularx}{\textwidth}{@{}p{3.2cm}X@{}}
\toprule
\textbf{Module} & \textbf{Responsabilité} \\
\midrule
\rowcolor{rowgray}
AuthModule & Login, register, refresh, logout (Argon2id, JWT, cookies) \\
AuthorizationModule & JwtAuthGuard, PermissionsGuard, isolation org \\
\rowcolor{rowgray}
AgentsModule & CRUD identités agents non-humaines \\
VendorsModule & Fournisseurs SaaS \\
\rowcolor{rowgray}
PoliciesModule & Politiques + évaluation via policy-engine \\
CredentialsModule & Coffre-fort, grants, métadonnées \\
\rowcolor{rowgray}
WorkflowsModule & Définitions de workflows (templates) \\
WorkflowRunsModule & Exécutions, état, annulation \\
\rowcolor{rowgray}
ApprovalsModule & Cycle de vie approbations humaines \\
QueueModule & Enqueue BullMQ (start, resume, cancel) \\
\rowcolor{rowgray}
InternalWorkerModule & Endpoints worker (events, screenshots, vault decrypt) \\
AuditModule & Événements audit + chaîne hash SHA-256 \\
\rowcolor{rowgray}
ReceiptsModule & Artefacts finaux de confiance \\
FilesModule & Métadonnées + stockage S3 \\
\bottomrule
\end{tabularx}
\end{table}

\textbf{Middleware global (ordre)~:}
\begin{enumerate}[leftmargin=1.8cm, itemsep=2pt]
  \item \texttt{RequestContextMiddleware} — user, org, request ID
  \item \texttt{SecurityHeadersMiddleware}
  \item \texttt{RateLimitMiddleware} — Redis TTL en production
  \item \texttt{RequestLoggingMiddleware}
\end{enumerate}

\subsection{Architecture physique et déploiement}

\begin{figure}[H]
\centering
\diagramespace[6.5cm]{Docker Compose — postgres:17 (5432), redis:8 (6379),
  minio (9000/9001), mailpit (1025/8025), api (3001), worker, web (3000),
  vendor-sandbox (4202), migrate — profils infra et app}
\caption{Architecture de déploiement Docker Compose d'AegisWeb}
\label{fig:archi_deploy}
\end{figure}

\begin{table}[H]
\centering
\caption{Services conteneurisés et ports locaux}
\label{tab:docker}
\begin{tabularx}{\textwidth}{@{}p{2.2cm}p{3.2cm}X@{}}
\toprule
\textbf{Service} & \textbf{Image/Port}     & \textbf{Rôle} \\
\midrule
\rowcolor{rowgray}
\texttt{postgres} & postgres:17 / 5432   & Base \texttt{agentpass} \\
\texttt{redis}    & redis:8 / 6379       & BullMQ + rate limit \\
\rowcolor{rowgray}
\texttt{minio}    & minio / 9000, 9001   & Stockage S3-compatible \\
\texttt{mailpit}  & axllent/mailpit      & SMTP 1025, UI 8025 \\
\rowcolor{rowgray}
\texttt{api}      & NestJS / 3001        & API REST + Swagger \\
\texttt{worker}  & NestJS + Playwright  & Consommation jobs \\
\rowcolor{rowgray}
\texttt{web}      & Next.js / 3000       & Dashboard \\
\texttt{vendor-sandbox} & NestJS / 4202 & Portail SaaS factice \\
\bottomrule
\end{tabularx}
\end{table}

\section{Modélisation dynamique}

\subsection{Diagramme de séquence — Authentification}

\begin{figure}[H]
\centering
\diagramespace[7cm]{Séquence Auth — Client $\rightarrow$ POST /auth/login
  $\rightarrow$ AuthService (Argon2id verify) $\rightarrow$ JWT access
  (body) + refresh (cookie HttpOnly) $\rightarrow$ cookie
  aegisweb\_session (middleware) — branche 401 $\rightarrow$ /auth/refresh}
\caption{Diagramme de séquence — Authentification JWT + refresh}
\label{fig:seq_auth}
\end{figure}

\subsection{Diagramme de séquence — Exécution workflow avec approbation}

\begin{figure}[H]
\centering
\diagramespace[8cm]{Séquence Workflow — Owner $\rightarrow$ POST
  /workflow-runs $\rightarrow$ QueueService.enqueue $\rightarrow$ Worker
  $\rightarrow$ Playwright $\rightarrow$ PolicyEngine.evaluate $\rightarrow$
  [require\_approval] $\rightarrow$ ApprovalRequest $\rightarrow$ Email
  $\rightarrow$ Approver $\rightarrow$ POST /approvals/:id/approve
  $\rightarrow$ resume queue $\rightarrow$ Receipt}
\caption{Diagramme de séquence — Workflow downgrade avec approbation}
\label{fig:seq_workflow}
\end{figure}

\subsection{Diagramme d'activités — Évaluation du moteur de politiques}

\begin{figure}[H]
\centering
\diagramespace[6.5cm]{Activités Policy Engine — ActionAttempt reçue
  $\rightarrow$ Agent actif ? $\rightarrow$ Domaine bloqué ? $\rightarrow$
  Allowlist OK ? $\rightarrow$ Action interdite ? $\rightarrow$ Seuil
  montant ? $\rightarrow$ [allow | deny | require\_approval | pause\_agent]
  + RiskLevel + signaux}
\caption{Diagramme d'activités — Pipeline d'évaluation des politiques}
\label{fig:act_policy}
\end{figure}

\section{Modélisation statique}

\subsection{Modèle de données Prisma — Vue d'ensemble}

Le schéma (\texttt{prisma/schema.prisma}) comprend \textbf{17 modèles}
principaux et de nombreux enums typés.

\begin{table}[H]
\centering
\caption{Domaines fonctionnels et entités Prisma}
\label{tab:domaines_bd}
\begin{tabularx}{\textwidth}{@{}p{2.4cm}X@{}}
\toprule
\textbf{Domaine}  & \textbf{Modèles / Entités} \\
\midrule
\rowcolor{rowgray}
Identité          & \texttt{Organization}, \texttt{User}, \texttt{RefreshToken} \\
Agents            & \texttt{Agent} \\
\rowcolor{rowgray}
Fournisseurs      & \texttt{Vendor} \\
Politiques        & \texttt{Policy} \\
\rowcolor{rowgray}
Credentials       & \texttt{Credential}, \texttt{CredentialAgentGrant} \\
Workflows         & \texttt{Workflow}, \texttt{WorkflowRun}, \texttt{ActionAttempt} \\
\rowcolor{rowgray}
Approbations      & \texttt{ApprovalRequest} \\
Audit \& Preuves  & \texttt{AuditEvent}, \texttt{File}, \texttt{Receipt} \\
\bottomrule
\end{tabularx}
\end{table}

\subsection{Enums métier principaux}

\begin{table}[H]
\centering
\caption{Énumérations métier clés}
\label{tab:enums}
\small
\begin{tabularx}{\textwidth}{@{}p{3.2cm}X@{}}
\toprule
\textbf{Enum} & \textbf{Valeurs principales} \\
\midrule
\rowcolor{rowgray}
\texttt{UserRole} & owner, admin, approver, auditor, developer \\
\texttt{WorkflowTemplate} & vendor\_invoice\_download, saas\_renewal\_check,
  plan\_downgrade\_request \\
\rowcolor{rowgray}
\texttt{WorkflowRunStatus} & queued, running, waiting\_for\_approval,
  completed, failed, canceled, denied \\
\texttt{PolicyDecision} & allow, deny, require\_approval,
  require\_step\_up\_auth, pause\_agent \\
\rowcolor{rowgray}
\texttt{RiskLevel} & low, medium, high, critical \\
\texttt{ActionType} & open\_page, read\_page, fill\_form, click\_button,
  download\_file, change\_plan, cancel\_subscription, \ldots \\
\bottomrule
\end{tabularx}
\end{table}

\subsection{Dictionnaire de données — Entité \texttt{WorkflowRun}}

\begin{table}[H]
\centering
\caption{Dictionnaire de données — Modèle \texttt{WorkflowRun}}
\label{tab:dico_run}
\small
\begin{tabularx}{\textwidth}{@{}p{3.6cm}p{2.8cm}X@{}}
\toprule
\textbf{Attribut} & \textbf{Type} & \textbf{Description} \\
\midrule
\texttt{id} & UUID & Clé primaire \\
\rowcolor{rowgray}
\texttt{organizationId} & UUID (FK) & Isolation multi-tenant \\
\texttt{workflowId} & UUID (FK) & Définition parente \\
\rowcolor{rowgray}
\texttt{agentId} & UUID (FK) & Agent exécutant \\
\texttt{status} & WorkflowRunStatus & État courant du run \\
\rowcolor{rowgray}
\texttt{stateJson} & JSON & État interne sérialisé (étapes, contexte) \\
\texttt{requestedByUserId} & UUID & Utilisateur ayant lancé le run \\
\rowcolor{rowgray}
\texttt{startedAt} / \texttt{completedAt} & DateTime & Horodatages \\
\texttt{failureReason} & String? & Raison d'échec si failed \\
\bottomrule
\end{tabularx}
\end{table}

\subsection{Dictionnaire de données — Entité \texttt{AuditEvent}}

\begin{table}[H]
\centering
\caption{Dictionnaire de données — Modèle \texttt{AuditEvent}}
\label{tab:dico_audit}
\small
\begin{tabularx}{\textwidth}{@{}p{3.6cm}p{2.8cm}X@{}}
\toprule
\textbf{Attribut} & \textbf{Type} & \textbf{Description} \\
\midrule
\texttt{id} & UUID & Clé primaire \\
\rowcolor{rowgray}
\texttt{organizationId} & UUID (FK) & Tenant \\
\texttt{eventType} & AuditEventType & Type typé (50+ événements) \\
\rowcolor{rowgray}
\texttt{actorType} & AuditActorType & user, agent, worker, system \\
\texttt{actorId} & String? & Identifiant de l'acteur \\
\rowcolor{rowgray}
\texttt{prevHash} & String & Hash événement précédent (chaîne) \\
\texttt{eventHash} & String & SHA-256 de l'événement courant \\
\rowcolor{rowgray}
\texttt{payloadJson} & JSON & Données redactées (secrets masqués) \\
\texttt{createdAt} & DateTime & Horodatage append-only \\
\bottomrule
\end{tabularx}
\end{table}

\subsection{Système de files BullMQ}

\begin{table}[H]
\centering
\caption{Files BullMQ et types de jobs}
\label{tab:queues}
\begin{tabularx}{\textwidth}{@{}p{3.2cm}Xp{3.2cm}@{}}
\toprule
\textbf{Queue} & \textbf{Usage} & \textbf{ID job} \\
\midrule
\rowcolor{rowgray}
\texttt{workflow-runs} & Démarrage de workflows &
  \texttt{start-\{runId\}} \\
\texttt{workflow-resume} & Reprise après approbation &
  \texttt{resume-\{runId\}-\{approvalId\}} \\
\rowcolor{rowgray}
\texttt{workflow-maintenance} & Signaux d'annulation &
  \texttt{cancel-\{runId\}} \\
\bottomrule
\end{tabularx}
\end{table}

\subsection{RBAC — Matrice des permissions}

\begin{table}[H]
\centering
\caption{Rôles utilisateur et permissions principales}
\label{tab:rbac}
\small
\begin{tabularx}{\textwidth}{@{}p{2cm}X@{}}
\toprule
\textbf{Rôle} & \textbf{Permissions} \\
\midrule
\rowcolor{rowgray}
owner / admin & Toutes (28 permissions) \\
approver & Lecture + \texttt{approval:approve} \\
\rowcolor{rowgray}
auditor & Lecture + \texttt{audit:read} + \texttt{file:read} \\
developer & Lecture + \texttt{workflow:run} + \texttt{workflow:cancel} \\
\bottomrule
\end{tabularx}
\end{table}

\subsection{Schéma entité-relation}

\begin{figure}[H]
\centering
\diagramespace[5.5cm]{ER simplifié — Organization (1)--(*) User, Agent,
  Vendor, Policy, Workflow — Agent (*)--(*) Credential via Grant —
  Workflow (1)--(*) WorkflowRun (1)--(*) ActionAttempt, ApprovalRequest,
  Receipt — WorkflowRun (1)--(*) AuditEvent, File}
\caption{Schéma entité-relation simplifié — base de données AegisWeb}
\label{fig:er}
\end{figure}

\conclusionchap{La conception d'AegisWeb révèle une architecture modulaire
et sécurisée~: monorepo TypeScript, séparation control plane / worker,
moteur de politiques pur, vault chiffré et modèle de données Prisma riche
en 17 tables. Ces fondements guident notre implémentation, décrite dans
le chapitre suivant.}

% =============================================================
% CHAPITRE 4 — RÉALISATION
% =============================================================
\chapter{Réalisation}

\begin{chapintro}
Ce chapitre présente la réalisation technique de la plateforme AegisWeb.
Après la description de l'environnement et des technologies, nous exposons
les livrables des phases principales, illustrons les interfaces du
dashboard et présentons les résultats des tests de validation.
\end{chapintro}

\section{Environnement de développement}

\subsection{Environnement matériel}

\begin{table}[H]
\centering
\caption{Spécifications de l'environnement matériel}
\label{tab:materiel}
\begin{tabularx}{\textwidth}{@{}p{4cm}X@{}}
\toprule
\textbf{Composant}       & \textbf{Spécification} \\
\midrule
\rowcolor{rowgray}
Processeur               & Intel Core i7 / AMD Ryzen 7 (multi-cœurs) \\
Mémoire RAM              & 16~Go minimum (32~Go recommandé pour Playwright) \\
\rowcolor{rowgray}
Stockage                 & SSD NVMe 512~Go \\
Système d'exploitation   & Linux (Ubuntu) / Windows 11 + WSL2 \\
\rowcolor{rowgray}
Connexion réseau         & Internet pour installation dépendances npm \\
\bottomrule
\end{tabularx}
\end{table}

\subsection{Environnement logiciel}

\begin{table}[H]
\centering
\caption{Environnement logiciel et outils de développement}
\label{tab:logiciel}
\begin{tabularx}{\textwidth}{@{}p{4cm}X@{}}
\toprule
\textbf{Catégorie}       & \textbf{Outil / Version} \\
\midrule
\rowcolor{rowgray}
Runtime                  & Node.js $\geq$ 22 \\
Gestionnaire paquets     & pnpm $\geq$ 10 (workspaces) \\
\rowcolor{rowgray}
Éditeur                  & Visual Studio Code / Cursor \\
Langage                  & TypeScript 5.7+ \\
\rowcolor{rowgray}
Conteneurisation         & Docker Compose v2 \\
Base de données          & PostgreSQL 17, Prisma 7.8 \\
\rowcolor{rowgray}
Tests                    & Vitest 2.1 + Supertest \\
Lint                     & ESLint 9 \\
\rowcolor{rowgray}
Navigateur automatisé    & Playwright Chromium (worker) \\
Client API               & Swagger (\texttt{/docs}), Bruno/Postman \\
\bottomrule
\end{tabularx}
\end{table}

\section{Technologies utilisées}

\subsection{Backend — API NestJS}

\begin{table}[H]
\centering
\caption{Technologies du backend AegisWeb}
\label{tab:tech_backend}
\begin{tabularx}{\textwidth}{@{}p{3.6cm}p{2.4cm}X@{}}
\toprule
\textbf{Technologie}      & \textbf{Version} & \textbf{Rôle} \\
\midrule
\rowcolor{rowgray}
Node.js                   & 22+              & Runtime JavaScript serveur \\
NestJS                    & 11               & Framework modulaire (DI, guards) \\
\rowcolor{rowgray}
Prisma                    & 7.8              & ORM, migrations, client typé \\
PostgreSQL                & 17               & Base relationnelle \\
\rowcolor{rowgray}
Redis + BullMQ            & 8 / 5            & Files de workflows \\
AWS SDK (S3)              & 3.x              & Stockage MinIO/S3 \\
\rowcolor{rowgray}
Argon2                    & —                & Hachage mots de passe \\
jsonwebtoken              & —                & JWT access tokens \\
\rowcolor{rowgray}
Zod                       & 3.24+            & Validation schémas \\
Vitest + Supertest        & 2.1+             & Tests intégration API \\
\bottomrule
\end{tabularx}
\end{table}

\subsection{Worker et runtime navigateur}

\begin{table}[H]
\centering
\caption{Technologies du worker et browser-runtime}
\label{tab:tech_worker}
\begin{tabularx}{\textwidth}{@{}p{3.6cm}p{2.4cm}X@{}}
\toprule
\textbf{Technologie}      & \textbf{Version}  & \textbf{Rôle} \\
\midrule
\rowcolor{rowgray}
NestJS application context & 11               & Boot worker sans HTTP public \\
BullMQ Worker             & 5                 & Consommation jobs (concurrency=1) \\
\rowcolor{rowgray}
Playwright                & 1.60              & Automatisation Chromium \\
libs/browser-runtime      & —                 & Allowlist, screenshots, downloads \\
\rowcolor{rowgray}
libs/vault                & —                 & Déchiffrement credential scopé run \\
HMAC worker tokens        & —                 & Auth endpoints internes (6h TTL) \\
\bottomrule
\end{tabularx}
\end{table}

\subsection{Frontend — Dashboard Next.js}

\begin{table}[H]
\centering
\caption{Technologies du frontend AegisWeb}
\label{tab:tech_frontend}
\begin{tabularx}{\textwidth}{@{}p{3.6cm}p{2.4cm}X@{}}
\toprule
\textbf{Technologie}      & \textbf{Version}  & \textbf{Rôle} \\
\midrule
\rowcolor{rowgray}
Next.js                   & 16                & App Router, SSR, middleware \\
React                     & 19                & Composants UI \\
\rowcolor{rowgray}
Tailwind CSS              & 4                 & Styles utilitaires \\
Radix UI / shadcn         & —                 & Primitives accessibles (\texttt{components/ui}) \\
\rowcolor{rowgray}
TanStack React Query      & 5                 & Cache serveur, polling runs/approvals \\
react-hook-form + Zod     & —                 & Formulaires typés \\
\rowcolor{rowgray}
Sonner                    & —                 & Notifications toast \\
Iconoir + Lucide          & —                 & Iconographie \\
\bottomrule
\end{tabularx}
\end{table}

\section{Phases de réalisation et livrables}

\subsection{Phases 0--6 — Fondation, auth et audit}

\begin{table}[H]
\centering
\caption{Phases 0 à 6 — Fondations et sécurité transversale}
\label{tab:phases_0_6}
\begin{tabularx}{\textwidth}{@{}cXp{2cm}c@{}}
\toprule
\textbf{Ph.} & \textbf{Livrable} & \textbf{Tests} & \textbf{Statut} \\
\midrule
0 & Boot API, worker, sandbox, libs & phase0.spec.ts & \coche \\
\rowcolor{rowgray}
1 & Prisma schema + seed Northstar Labs & phase1-database & \coche \\
2 & Lib domain (enums, permissions) & phase2-domain & \coche \\
\rowcolor{rowgray}
3 & Config, logging, health, erreurs & phase3-api-foundation & \coche \\
4 & AuthModule Argon2id + JWT + refresh & phase4-auth & \coche \\
\rowcolor{rowgray}
5 & RBAC + isolation organisation & phase5-authorization & \coche \\
6 & AuditModule hash chain SHA-256 & phase6-audit & \coche \\
\bottomrule
\end{tabularx}
\end{table}

\textbf{Réalisations clés~:}
\begin{itemize}[leftmargin=1.8cm, itemsep=3pt]
  \item Stack Docker Compose (Postgres, Redis, MinIO, Mailpit).
  \item Format réponse API unifié~: \texttt{\{ data \}} / \texttt{\{ error \}}.
  \item Refresh token atomique (migration \texttt{atomic\_refresh\_token\_hash}).
  \item Chaîne d'audit~: chaque \texttt{AuditEvent} référence \texttt{prevHash}.
\end{itemize}

\subsection{Phases 7--18 — Ressources métier et queue}

\begin{table}[H]
\centering
\caption{Phases 7 à 18 — Modules CRUD et orchestration}
\label{tab:phases_7_18}
\begin{tabularx}{\textwidth}{@{}cXp{2cm}c@{}}
\toprule
\textbf{Ph.} & \textbf{Livrable} & \textbf{Tests} & \textbf{Statut} \\
\midrule
\rowcolor{rowgray}
7 & FilesModule S3/MinIO & phase7-files & \coche \\
8 & Organization + Users & phase8-organization-users & \coche \\
\rowcolor{rowgray}
9--10 & Agents + Vendors & phase9, phase10 & \coche \\
11--12 & Policy Engine + PoliciesModule & phase11, phase12 & \coche \\
\rowcolor{rowgray}
13 & Vault AES-GCM + Credentials & phase13-credentials-vault & \coche \\
14--15 & Workflows + WorkflowRuns & phase14, phase15 & \coche \\
\rowcolor{rowgray}
16--17 & ActionAttempts + Approvals & phase16, phase17 & \coche \\
18 & QueueModule BullMQ & phase18-queue & \coche \\
\bottomrule
\end{tabularx}
\end{table}

\subsection{Phases 19--29 — Worker, workflows et finalisation}

\begin{table}[H]
\centering
\caption{Phases 19 à 29 — Worker, browser et workflows SaaS}
\label{tab:phases_19_29}
\begin{tabularx}{\textwidth}{@{}cXp{2cm}c@{}}
\toprule
\textbf{Ph.} & \textbf{Livrable} & \textbf{Tests} & \textbf{Statut} \\
\midrule
\rowcolor{rowgray}
19--20 & Worker foundation + browser-runtime & phase19, phase20 & \coche \\
21--22 & Vendor sandbox + Connector & phase21, phase22 & \coche \\
\rowcolor{rowgray}
23 & InternalWorkerModule & phase23-internal-worker & \coche \\
24--26 & 3 workflows SaaS + approbation & phase24--26 & \coche \\
\rowcolor{rowgray}
27--28 & Receipts + Notifications email & phase27, phase28 & \coche \\
29 & OpenAPI / Swagger & phase29-api-docs & \coche \\
\bottomrule
\end{tabularx}
\end{table}

\textbf{Workflows implémentés~:}
\begin{itemize}[leftmargin=1.8cm, itemsep=3pt]
  \item \texttt{vendor\_invoice\_download}~: login $\rightarrow$ téléchargement
    facture PDF $\rightarrow$ screenshot $\rightarrow$ receipt.
  \item \texttt{saas\_renewal\_check}~: login $\rightarrow$ lecture données
    renouvellement $\rightarrow$ screenshot $\rightarrow$ receipt.
  \item \texttt{plan\_downgrade\_request}~: préparation changement plan
    $\rightarrow$ pause approbation $\rightarrow$ reprise $\rightarrow$ soumission.
\end{itemize}

\section{Extraits de code significatifs}

\subsection{Permissions RBAC partagées}

\begin{lstlisting}[style=tsstyle,
  caption={libs/domain — Définition des permissions et rôles},
  label={lst:permissions}]
export const Permission = {
  AgentCreate: 'agent:create',
  AgentRead: 'agent:read',
  CredentialGrant: 'credential:grant',
  WorkflowRun: 'workflow:run',
  ApprovalApprove: 'approval:approve',
  AuditRead: 'audit:read',
  // ... 28 permissions au total
} as const;

export const ROLE_PERMISSIONS: Record<UserRole, readonly Permission[]> = {
  [UserRole.Owner]: PERMISSIONS,
  [UserRole.Approver]: [
    Permission.ApprovalRead,
    Permission.ApprovalApprove,
    Permission.ReceiptRead,
    // ...
  ],
  [UserRole.Developer]: [
    ...readPermissions,
    Permission.WorkflowRun,
    Permission.WorkflowCancel,
  ],
};
\end{lstlisting}

\subsection{Middleware de sécurité frontend}

\begin{lstlisting}[style=tsstyle,
  caption={apps/web/middleware.ts — Protection routes /app et CSP},
  label={lst:middleware}]
export function middleware(request: NextRequest) {
  const nonce = generateNonce();
  const response = NextResponse.next();

  // Routes dashboard protegees
  if (request.nextUrl.pathname.startsWith('/app')) {
    const session = request.cookies.get('aegisweb_session');
    if (!session) {
      return NextResponse.redirect(new URL('/login', request.url));
    }
  }

  // CSP avec nonce (pas de unsafe-eval)
  response.headers.set('Content-Security-Policy', buildCsp(nonce));
  response.headers.set('X-Frame-Options', 'DENY');
  return response;
}
\end{lstlisting}

\subsection{Enqueue idempotent BullMQ}

\begin{lstlisting}[style=tsstyle,
  caption={Queue — Jobs idempotents par runId},
  label={lst:queue}]
async enqueueWorkflowStart(runId: string, payload: StartPayload) {
  await this.workflowRunsQueue.add(
    'start-workflow-run',
    payload,
    {
      jobId: `start-${runId}`,
      attempts: 3,
      backoff: { type: 'exponential', delay: 2000 },
    },
  );
}
\end{lstlisting}

\section{Interfaces du dashboard}

Le dashboard Next.js (\texttt{apps/web}) suit une direction UX
\textbf{evidence-first}~: chaque écran opérationnel met en avant les
preuves (screenshots, timeline, audit drawer, receipts).

\subsection{Structure des routes}

\begin{table}[H]
\centering
\caption{Routes principales du dashboard (\texttt{/app/*})}
\label{tab:routes}
\small
\begin{tabularx}{\textwidth}{@{}p{4.2cm}X@{}}
\toprule
\textbf{Route} & \textbf{Fonction} \\
\midrule
\rowcolor{rowgray}
\texttt{/app/home} & Tableau de bord opérationnel \\
\texttt{/app/agents} & Gestion identités agents \\
\rowcolor{rowgray}
\texttt{/app/vendors} & Fournisseurs SaaS \\
\texttt{/app/credentials} & Coffre-fort (sans secrets en clair) \\
\rowcolor{rowgray}
\texttt{/app/policies} & Politiques de gouvernance \\
\texttt{/app/workflows} & Définitions de workflows \\
\rowcolor{rowgray}
\texttt{/app/runs} & Exécutions avec polling état \\
\texttt{/app/approvals} & File d'approbation humaine \\
\rowcolor{rowgray}
\texttt{/app/receipts} & Reçus de confiance \\
\texttt{/app/audit} & Journal d'audit chaîné \\
\bottomrule
\end{tabularx}
\end{table}

\subsection{Captures d'interface}

\begin{figure}[H]
\centering
\begin{minipage}[t]{0.47\textwidth}
  \diagramespace[7.5cm]{Capture dashboard — Page d'accueil /app/home
    (métriques runs, approvals en attente, activité récente)}
  \caption{Tableau de bord opérationnel}
  \label{fig:ui_home}
\end{minipage}
\hfill
\begin{minipage}[t]{0.47\textwidth}
  \diagramespace[7.5cm]{Capture dashboard — Détail run avec timeline
    evidence, screenshots, décisions policy, statut approval}
  \caption{Détail d'exécution workflow — panneau preuves}
  \label{fig:ui_run}
\end{minipage}
\end{figure}

\begin{figure}[H]
\centering
\begin{minipage}[t]{0.47\textwidth}
  \diagramespace[7.5cm]{Capture dashboard — Formulaire création agent
    (identifier, purpose, status)}
  \caption{Gestion des agents IA}
  \label{fig:ui_agents}
\end{minipage}
\hfill
\begin{minipage}[t]{0.47\textwidth}
  \diagramespace[7.5cm]{Capture dashboard — Écran approbation avec
    screenshots billing, boutons Approve/Reject, contexte risque}
  \caption{Interface d'approbation humaine}
  \label{fig:ui_approval}
\end{minipage}
\end{figure}

\begin{figure}[H]
\centering

\includegraphics[width=0.35\textwidth]{assets/aegisweb_primary_logo.png}
\caption{Logo principal AegisWeb — identité visuelle de la plateforme}
\label{fig:logo_primary}
\end{figure}

\section{Données de démonstration}

Après \texttt{pnpm db:reset}, l'environnement local contient~:

\begin{table}[H]
\centering
\caption{Données seed — Organisation Northstar Labs}
\label{tab:seed}
\small
\begin{tabularx}{\textwidth}{@{}lX@{}}
\toprule
\textbf{Élément} & \textbf{Valeur} \\
\midrule
\rowcolor{rowgray}
Organisation & Northstar Labs (\texttt{northstarlabs.dev}) \\
Utilisateurs & founder (owner), finance (approver), auditor, dev
  — mot de passe \texttt{Password123!} \\
\rowcolor{rowgray}
Agents & Procurement Bot, Invoice Collector, Risky Admin Bot,
  Legacy Spend Bot \\
Vendors & Acme Analytics, Nimbus Docs, Atlas CRM, PayrollPro \\
\rowcolor{rowgray}
Workflows & Acme Invoice Download, Acme Renewal Check,
  Acme Downgrade Request, Atlas Renewal Check \\
\bottomrule
\end{tabularx}
\end{table}

\section{Tests et validation}

\subsection{Stratégie de tests}

\begin{itemize}[leftmargin=1.8cm, itemsep=3pt]
  \item \textbf{33 fichiers Vitest} dans \texttt{tests/} (phases 0--29 +
    acceptance + sécurité + dashboard client).
  \item Environnement \texttt{node} pour backend/worker ;
    \texttt{jsdom} pour composants web.
  \item Pool forks single-fork pour tests partageant services locaux.
  \item Scripts complémentaires~: \texttt{pnpm smoke}, \texttt{pnpm e2e:happy},
    \texttt{pnpm qa:click-path}, \texttt{pnpm prod:check}.
\end{itemize}

\subsection{Résultats d'acceptation MVP}

\begin{table}[H]
\centering
\caption{Résultats de la suite d'acceptation MVP backend}
\label{tab:acceptance}
\begin{tabularx}{\textwidth}{@{}lX@{}}
\toprule
\textbf{Métrique} & \textbf{Valeur documentée} \\
\midrule
\rowcolor{rowgray}
Fichiers de tests & 31 (+ security + web client) \\
Tests passés & 195 \\
\rowcolor{rowgray}
Fichier acceptance & \texttt{mvp-backend-acceptance.spec.ts} \\
Régression sécurité & \texttt{security-coverage.spec.ts} \\
\bottomrule
\end{tabularx}
\end{table}

\subsection{Tests fonctionnels représentatifs}

\begin{table}[H]
\centering
\caption{Tests fonctionnels backend représentatifs}
\label{tab:tests_backend}
\small
\begin{tabularx}{\textwidth}{@{}p{0.9cm}Xp{3.6cm}c@{}}
\toprule
\textbf{ID} & \textbf{Scénario} & \textbf{Résultat attendu} & \textbf{Statut} \\
\midrule
\rowcolor{rowgray}
T01 & Login owner Northstar Labs & JWT + cookie refresh & \coche \\
T02 & Isolation cross-org & HTTP 404/403 & \coche \\
\rowcolor{rowgray}
T03 & Création agent + grant credential & Grant actif en BDD & \coche \\
T04 & Policy deny action interdite & ActionAttempt denied & \coche \\
\rowcolor{rowgray}
T05 & Workflow invoice download E2E & Run completed + receipt & \coche \\
T06 & Downgrade + approbation & Pause $\rightarrow$ resume $\rightarrow$ receipt & \coche \\
\rowcolor{rowgray}
T07 & Audit hash chain integrity & prevHash cohérent & \coche \\
T08 & Vault decrypt scopé run & Secret déchiffré worker only & \coche \\
\rowcolor{rowgray}
T09 & Browser allowlist violation & Action bloquée & \coche \\
T10 & Refresh token révoqué & HTTP 401 & \coche \\
\bottomrule
\end{tabularx}
\end{table}

\subsection{Score audit sécurité}

\begin{table}[H]
\centering
\caption{Scores d'audit sécurité (docs/SECURITY\_AUDIT.md)}
\label{tab:security_audit}
\begin{tabularx}{\textwidth}{@{}p{4.5cm}C{2.2cm}@{}}
\toprule
\textbf{Domaine} & \textbf{Score /5} \\
\midrule
\rowcolor{rowgray}
Credential vault & 4,0 \\
Auth / sessions & 3,8 \\
\rowcolor{rowgray}
Autorisation / tenant & 4,0 \\
Worker / internal API & 3,8 \\
\rowcolor{rowgray}
Browser runtime & 3,8 \\
Evidence / files & 3,3 \\
\rowcolor{rowgray}
Frontend & 3,8 \\
\textbf{Pilot readiness global} & \textbf{3,8} \\
\bottomrule
\end{tabularx}
\end{table}

\subsection{Métriques de performance}

\begin{table}[H]
\centering
\caption{Métriques mesurées en environnement de développement local}
\label{tab:perfs}
\begin{tabularx}{\textwidth}{@{}p{5cm}Xp{2.6cm}@{}}
\toprule
\textbf{Métrique} & \textbf{Description} & \textbf{Valeur} \\
\midrule
\rowcolor{rowgray}
Readiness API & GET \texttt{/health/ready} & Postgres + Redis OK \\
Latence CRUD & GET agents/vendors & $<$ 100 ms typique \\
\rowcolor{rowgray}
Workflow E2E & Downgrade avec approbation & 30--90 s (Playwright) \\
Build Next.js & \texttt{pnpm build:web} & $\approx$ 2--4 min \\
\rowcolor{rowgray}
Suite Vitest complète & \texttt{pnpm test} & $\approx$ 3--8 min \\
\bottomrule
\end{tabularx}
\end{table}

\conclusionchap{Ce chapitre a présenté la réalisation complète d'AegisWeb
sur 30 phases incrémentales, couvrant l'API NestJS, le worker Playwright,
les bibliothèques de sécurité et le dashboard Next.js. L'intégration
PostgreSQL, Redis, BullMQ, MinIO et Prisma a produit un MVP local
fonctionnel, validé par 195 tests Vitest et un audit sécurité documenté
(3,8/5 pilot readiness).}

% =============================================================
% CONCLUSION GÉNÉRALE
% =============================================================
\clearpage
\chapter*{Conclusion Générale}
\addcontentsline{toc}{chapter}{Conclusion Générale}
\markboth{Conclusion Générale}{Conclusion Générale}

\begin{chapintro}
Ce projet de fin d'études nous a offert l'opportunité de concevoir et
de développer, de bout en bout, une plateforme de gouvernance pour agents
IA agissant sur le web — un domaine émergent à l'intersection de la
cybersécurité, de l'automatisation et de l'intelligence artificielle.
\end{chapintro}

\medskip
\noindent\textbf{\large Bilan du projet}

\medskip
Au terme de ce travail, nous avons conçu et développé \textbf{AegisWeb},
une plateforme de confiance pour agents IA web répondant au principe
\textit{«~Permission before autonomy~»}. L'ensemble des objectifs initiaux
a été atteint :

\begin{itemize}[leftmargin=1.8cm, itemsep=4pt]
  \item Un \textbf{monorepo TypeScript} structuré (4 applications, 7
    bibliothèques partagées) avec pnpm workspaces.
  \item Une \textbf{API NestJS} modulaire (20+ modules) avec auth Argon2id,
    RBAC à 5 rôles et isolation multi-tenant.
  \item Un \textbf{modèle de données Prisma} en 17 tables avec migrations
    versionnées et seed de démonstration riche.
  \item Un \textbf{coffre-fort AES-256-GCM} et un \textbf{moteur de
    politiques pur} évaluant risques et approbations.
  \item Un \textbf{worker BullMQ + Playwright} exécutant trois workflows
    SaaS procurement avec preuves et receipts.
  \item Un \textbf{dashboard Next.js 16} evidence-first pour la gestion
    opérationnelle et l'approbation humaine.
  \item Une \textbf{infrastructure Docker Compose} reproductible et une
    suite de \textbf{195 tests Vitest} validant le MVP.
\end{itemize}

Sur le plan personnel, ce projet nous a permis de consolider nos
compétences en architecture logicielle distribuée, en sécurité
applicative (vault, audit chaîné, CSP), en développement fullstack
TypeScript (NestJS + Next.js) et en automatisation navigateur contrôlée.
Il nous a également sensibilisés aux enjeux de gouvernance des agents IA
dans un contexte entreprise.

\medskip
\noindent\textbf{\large Difficultés rencontrées}

\medskip
Ce projet n'a pas été exempt de difficultés :
\begin{itemize}[leftmargin=1.8cm, itemsep=3pt]
  \item L'\textbf{orchestration worker/API} (tokens HMAC, déchiffrement
    scopé, reprise après approbation) a exigé une conception rigoureuse
    des files BullMQ idempotentes.
  \item Le \textbf{runtime Playwright contrôlé} (allowlist, blocage
    réseau privé, masquage screenshots) a nécessité des garde-fous
    progressifs et des tests dédiés.
  \item La \textbf{cohérence frontend/backend} (polling runs, refresh
    token, middleware session) a demandé une couche données React Query
    structurée.
  \item Le \textbf{modèle de sécurité global} (redaction secrets, hash
    audit, isolation org) a requis des tests de régression continus.
\end{itemize}

\medskip
\noindent\textbf{\large Perspectives et améliorations futures}

\medskip
Plusieurs axes d'évolution sont envisagés :
\begin{itemize}[leftmargin=1.8cm, itemsep=3pt]
  \item \textbf{Passage production}~: checklist pilote
    (\texttt{SECURITY\_AUDIT.md}), Postgres/Redis managés, S3 privé,
    SMTP réel, déploiement Render documenté.
  \item \textbf{BFF HttpOnly}~: migration du stockage access token vers
    cookies HttpOnly pour réduire la surface XSS.
  \item \textbf{mTLS worker}~: authentification mutuelle entre worker et
    API interne en production.
  \item \textbf{Connecteurs vendors réels}~: extension au-delà du sandbox
    vers des portails SaaS réels avec step-up auth.
  \item \textbf{Observabilité}~: métriques Prometheus, tracing distribué
    des runs workflow.
  \item \textbf{Multi-région}~: réplication audit et stockage preuves pour
    conformité réglementaire.
\end{itemize}

% =============================================================
% BIBLIOGRAPHIE ET NETOGRAPHIE
% =============================================================
\clearpage
\chapter*{Bibliographie et Netographie}
\addcontentsline{toc}{chapter}{Bibliographie et Netographie}
\markboth{Bibliographie et Netographie}{Bibliographie et Netographie}

\section*{Bibliographie}

\begin{enumerate}[label={[\arabic*]}, leftmargin=1.8cm, itemsep=5pt]

  \item \textsc{Fowler}, M.,
    \textit{Patterns of Enterprise Application Architecture},
    Addison-Wesley, 2002.

  \item \textsc{Newman}, S.,
    \textit{Building Microservices},
    2\textsuperscript{e}~édition, O'Reilly Media, 2021.

  \item \textsc{Martin}, R.~C.,
    \textit{Clean Architecture: A Craftsman's Guide to Software
    Structure and Design}, Prentice Hall, 2017.

  \item \textsc{Richardson}, L. et \textsc{Ruby}, S.,
    \textit{RESTful Web Services}, O'Reilly Media, 2007.

  \item \textsc{Gamma}, E. et al.,
    \textit{Design Patterns: Elements of Reusable Object-Oriented
    Software}, Addison-Wesley, 1994.

  \item \textsc{Anderson}, R.~J.,
    \textit{Security Engineering},
    3\textsuperscript{e}~édition, Wiley, 2020.

  \item \textsc{Rubin}, K.~S.,
    \textit{Essential Scrum: A Practical Guide to the Most Popular
    Agile Process}, Addison-Wesley, 2012.

\end{enumerate}

\section*{Netographie}

\begin{enumerate}[label={[\arabic*]}, leftmargin=1.8cm, itemsep=5pt,
  resume]

  \item Documentation NestJS,
    \url{https://docs.nestjs.com/},
    [Consulté le : 15/05/2026].

  \item Documentation Next.js,
    \url{https://nextjs.org/docs},
    [Consulté le : 15/05/2026].

  \item Documentation Prisma,
    \url{https://www.prisma.io/docs},
    [Consulté le : 10/05/2026].

  \item Documentation BullMQ,
    \url{https://docs.bullmq.io/},
    [Consulté le : 12/05/2026].

  \item Documentation Playwright,
    \url{https://playwright.dev/docs/intro},
    [Consulté le : 18/05/2026].

  \item Documentation PostgreSQL 17,
    \url{https://www.postgresql.org/docs/17/},
    [Consulté le : 08/05/2026].

  \item OWASP Top Ten 2021,
    \url{https://owasp.org/www-project-top-ten/},
    [Consulté le : 20/05/2026].

  \item OWASP ASVS,
    \url{https://owasp.org/www-project-application-security-verification-standard/},
    [Consulté le : 20/05/2026].

  \item Documentation TanStack Query,
    \url{https://tanstack.com/query/latest},
    [Consulté le : 14/05/2026].

  \item Documentation Docker Compose,
    \url{https://docs.docker.com/compose/},
    [Consulté le : 05/05/2026].

  \item Documentation Redis,
    \url{https://redis.io/docs/},
    [Consulté le : 05/05/2026].

  \item RFC 9106 — Argon2 Memory-Hard Function,
    \url{https://www.rfc-editor.org/rfc/rfc9106},
    [Consulté le : 10/05/2026].

  \item NIST SP 800-38D — GCM Mode,
    \url{https://csrc.nist.gov/publications/detail/sp/800-38d/final},
    [Consulté le : 10/05/2026].

  \item Documentation Vitest,
    \url{https://vitest.dev/},
    [Consulté le : 22/05/2026].

  \item Dépôt AegisWeb — README et documentation technique,
    \url{https://github.com/},
    [Consulté le : 22/06/2026].

\end{enumerate}

% =============================================================
% ANNEXES
% =============================================================
\clearpage
\appendix

\chapter{Structure du projet — Arborescence des fichiers}

\begin{lstlisting}[
  basicstyle       = \ttfamily\scriptsize,
  frame            = single,
  numbers          = none,
  backgroundcolor  = \color{codebg},
  rulecolor        = \color{codeframe},
  caption          = {Arborescence du dépôt AegisWeb},
  label            = {lst:tree}
]
AegisWeb/
|-- apps/
|   |-- api/                    # API REST NestJS (control plane)
|   |   `-- src/
|   |       |-- main.ts
|   |       |-- app.module.ts
|   |       |-- auth/
|   |       |-- authorization/
|   |       |-- agents/
|   |       |-- vendors/
|   |       |-- policies/
|   |       |-- credentials/
|   |       |-- workflows/
|   |       |-- workflow-runs/
|   |       |-- approvals/
|   |       |-- receipts/
|   |       |-- audit/
|   |       |-- files/
|   |       |-- queue/
|   |       `-- internal-worker/
|   |-- worker/                 # Worker BullMQ + Playwright
|   |-- vendor-sandbox/         # Portail SaaS simulé
|   `-- web/                    # Dashboard Next.js 16
|       |-- app/                # App Router (/, /login, /app/*)
|       |-- components/         # UI, product, evidence, brand
|       |-- lib/                # auth, api, data-layer
|       `-- middleware.ts
|-- libs/
|   |-- domain/                 # Enums, permissions, types
|   |-- database/               # Prisma client + seed
|   |-- vault/                  # AES-256-GCM
|   |-- policy-engine/          # Moteur politiques pur
|   |-- browser-runtime/        # Playwright contrôlé
|   |-- audit/
|   `-- testing/
|-- prisma/
|   |-- schema.prisma
|   `-- migrations/
|-- infra/
|   `-- docker-compose.yml
|-- tests/                      # 33 specs Vitest (phase0--29)
|-- scripts/                    # demo-local, e2e-happy-path
`-- docs/                       # Plans phases, audits, runbooks
\end{lstlisting}

\chapter{Référence des endpoints API}

\begin{longtable}{@{}p{1.3cm}p{5.8cm}p{4cm}p{1.6cm}@{}}
\caption{Référence des endpoints API AegisWeb}
\label{tab:api_ref} \\
\toprule
\textbf{Méthode} & \textbf{Endpoint} & \textbf{Description} & \textbf{Auth} \\
\midrule
\endfirsthead
\toprule
\textbf{Méthode} & \textbf{Endpoint} & \textbf{Description} & \textbf{Auth} \\
\midrule
\endhead
\bottomrule
\endfoot
\bottomrule
\endlastfoot
\multicolumn{4}{@{}l}{\small\textit{Authentification (\texttt{/auth})}} \\[2pt]
POST & \texttt{/auth/register} & Inscription organisation & Non \\
POST & \texttt{/auth/login} & Connexion (JWT + refresh cookie) & Non \\
POST & \texttt{/auth/refresh} & Rafraîchir access token & Cookie \\
POST & \texttt{/auth/logout} & Révoquer refresh token & JWT \\
GET  & \texttt{/auth/me} & Profil utilisateur courant & JWT \\[4pt]
\multicolumn{4}{@{}l}{\small\textit{Santé (\texttt{/health})}} \\[2pt]
GET  & \texttt{/health} & Health check basique & Non \\
GET  & \texttt{/health/ready} & Readiness (Postgres, Redis) & Non \\[4pt]
\multicolumn{4}{@{}l}{\small\textit{Organisation (\texttt{/organization})}} \\[2pt]
GET  & \texttt{/organization} & Détails organisation & JWT \\
PATCH & \texttt{/organization} & Mise à jour organisation & JWT \\[4pt]
\multicolumn{4}{@{}l}{\small\textit{Utilisateurs (\texttt{/users})}} \\[2pt]
GET  & \texttt{/users} & Liste membres & JWT \\
POST & \texttt{/users/invite} & Inviter membre & JWT \\
PATCH & \texttt{/users/:id/role} & Modifier rôle & JWT \\[4pt]
\multicolumn{4}{@{}l}{\small\textit{Agents (\texttt{/agents})}} \\[2pt]
GET  & \texttt{/agents} & Liste agents & JWT \\
POST & \texttt{/agents} & Créer agent & JWT \\
GET  & \texttt{/agents/:id} & Détail agent & JWT \\
PATCH & \texttt{/agents/:id} & Modifier agent & JWT \\[4pt]
\multicolumn{4}{@{}l}{\small\textit{Vendors (\texttt{/vendors})}} \\[2pt]
GET/POST & \texttt{/vendors} & CRUD fournisseurs SaaS & JWT \\
GET/PATCH/DELETE & \texttt{/vendors/:id} & Opérations vendor & JWT \\[4pt]
\multicolumn{4}{@{}l}{\small\textit{Politiques (\texttt{/policies})}} \\[2pt]
GET/POST & \texttt{/policies} & CRUD politiques & JWT \\
POST & \texttt{/policies/:id/evaluate} & Évaluation politique & JWT \\[4pt]
\multicolumn{4}{@{}l}{\small\textit{Credentials (\texttt{/credentials})}} \\[2pt]
GET/POST & \texttt{/credentials} & CRUD credentials chiffrés & JWT \\
POST & \texttt{/credentials/:id/grants} & Accorder à un agent & JWT \\[4pt]
\multicolumn{4}{@{}l}{\small\textit{Workflows (\texttt{/workflows})}} \\[2pt]
GET/POST & \texttt{/workflows} & CRUD définitions workflow & JWT \\
GET/PATCH & \texttt{/workflows/:id} & Détail / mise à jour & JWT \\[4pt]
\multicolumn{4}{@{}l}{\small\textit{Exécutions (\texttt{/workflow-runs})}} \\[2pt]
GET/POST & \texttt{/workflow-runs} & Liste / lancer run & JWT \\
GET  & \texttt{/workflow-runs/:id} & Détail run + état & JWT \\
POST & \texttt{/workflow-runs/:id/cancel} & Annuler run & JWT \\
GET  & \texttt{/workflow-runs/:runId/action-attempts} & Tentatives & JWT \\[4pt]
\multicolumn{4}{@{}l}{\small\textit{Approbations (\texttt{/approvals})}} \\[2pt]
GET  & \texttt{/approvals} & File d'approbation & JWT \\
GET  & \texttt{/approvals/:id} & Détail approbation & JWT \\
POST & \texttt{/approvals/:id/approve} & Approuver & JWT \\
POST & \texttt{/approvals/:id/reject} & Rejeter & JWT \\[4pt]
\multicolumn{4}{@{}l}{\small\textit{Receipts, Audit, Files}} \\[2pt]
GET  & \texttt{/receipts} & Liste reçus de confiance & JWT \\
GET  & \texttt{/audit-events} & Journal d'audit & JWT \\
GET  & \texttt{/files/:id} & Métadonnées fichier & JWT \\[4pt]
\multicolumn{4}{@{}l}{\small\textit{Endpoints internes worker (\texttt{/internal})}} \\[2pt]
POST & \texttt{/internal/workers/runs/:runId/events} & Événements audit & Worker \\
POST & \texttt{/internal/workers/runs/:runId/screenshots} & Upload screenshot & Worker \\
POST & \texttt{/internal/workers/runs/:runId/complete} & Fin run réussie & Worker \\
POST & \texttt{/internal/workers/runs/:runId/fail} & Échec run & Worker \\
POST & \texttt{/internal/workers/runs/:runId/approval-requests} & Demande approbation & Worker \\
POST & \texttt{/internal/vault/credentials/:id/decrypt-for-run} & Déchiffrement scopé & Worker \\
\end{longtable}

\chapter{Variables d'environnement}

\begin{table}[H]
\centering
\caption{Variables d'environnement principales (\texttt{.env.example})}
\label{tab:env}
\small
\begin{tabularx}{\textwidth}{@{}p{4.8cm}X@{}}
\toprule
\textbf{Variable} & \textbf{Description} \\
\midrule
\rowcolor{rowgray}
\texttt{DATABASE\_URL} & Connexion PostgreSQL (\texttt{agentpass}) \\
\texttt{REDIS\_URL} & Connexion Redis pour BullMQ \\
\rowcolor{rowgray}
\texttt{JWT\_ACCESS\_SECRET} & Secret signature JWT access \\
\texttt{JWT\_REFRESH\_SECRET} & Secret refresh tokens \\
\rowcolor{rowgray}
\texttt{VAULT\_MASTER\_KEY} & Clé maître AES-256-GCM (44 chars base64) \\
\texttt{S3\_ENDPOINT} / \texttt{S3\_BUCKET} & Stockage MinIO/S3 \\
\rowcolor{rowgray}
\texttt{API\_PORT} / \texttt{API\_BASE\_URL} & API (défaut 3001) \\
\texttt{NEXT\_PUBLIC\_API\_URL} & URL API pour le frontend \\
\rowcolor{rowgray}
\texttt{WORKER\_INTERNAL\_TOKEN} & Token worker (dev) \\
\texttt{VENDOR\_SANDBOX\_URL} & URL sandbox (défaut 4202) \\
\rowcolor{rowgray}
\texttt{MAIL\_HOST} / \texttt{MAIL\_PORT} & SMTP Mailpit (1025) \\
\texttt{ENABLE\_OPENAPI} & Activer Swagger (\texttt{/docs}) \\
\rowcolor{rowgray}
\texttt{NEXT\_PUBLIC\_ENABLE\_DEMO\_MODE} & Mode démo dashboard \\
\texttt{API\_ALLOWED\_ORIGINS} & CORS origins autorisées \\
\bottomrule
\end{tabularx}
\end{table}

\chapter{Commandes de développement et déploiement}

\begin{table}[H]
\centering
\caption{Commandes pnpm principales}
\label{tab:commands}
\begin{tabularx}{\textwidth}{@{}p{4.2cm}X@{}}
\toprule
\textbf{Commande} & \textbf{Description} \\
\midrule
\rowcolor{rowgray}
\texttt{pnpm install} & Installation dépendances monorepo \\
\texttt{pnpm infra:up} & Docker Postgres, Redis, MinIO, Mailpit \\
\rowcolor{rowgray}
\texttt{pnpm db:reset} & Reset BDD + seed démo \\
\texttt{pnpm dev} & API + worker + web en parallèle \\
\rowcolor{rowgray}
\texttt{pnpm test} & Suite Vitest complète \\
\texttt{pnpm smoke} & Vérification services locaux \\
\rowcolor{rowgray}
\texttt{pnpm e2e:happy} & Parcours E2E heureux \\
\texttt{pnpm demo:local} & Démo locale one-command \\
\rowcolor{rowgray}
\texttt{pnpm prod:check} & typecheck + lint + build \\
\texttt{pnpm stack:hybrid:up} & Backend Docker + frontend local \\
\bottomrule
\end{tabularx}
\end{table}

\begin{table}[H]
\centering
\caption{URLs de services en développement local}
\label{tab:urls}
\begin{tabularx}{\textwidth}{@{}p{3.2cm}X@{}}
\toprule
\textbf{Service} & \textbf{URL} \\
\midrule
\rowcolor{rowgray}
Dashboard & \url{http://localhost:3000} \\
API & \url{http://localhost:3001} \\
\rowcolor{rowgray}
Readiness & \url{http://localhost:3001/health/ready} \\
Swagger & \url{http://localhost:3001/docs} \\
\rowcolor{rowgray}
Vendor sandbox & \url{http://localhost:4202} \\
Mailpit & \url{http://localhost:8025} \\
\rowcolor{rowgray}
MinIO console & \url{http://localhost:9001} \\
\bottomrule
\end{tabularx}
\end{table}


\end{document}
